% !TEX TS-program = pdflatexmk
\documentclass{article}

\usepackage[eandd,final]{neurips_2026}

\usepackage[utf8]{inputenc}
\usepackage[T1]{fontenc}
\usepackage{hyperref}
\usepackage{url}
\usepackage{booktabs}
\usepackage{amsfonts}
\usepackage{amssymb}
\usepackage{nicefrac}
\usepackage{microtype}
\usepackage{graphicx}
\usepackage{subcaption}
\usepackage{amsmath}
\usepackage{multirow}
\usepackage{xcolor}
\usepackage{colortbl}
\usepackage{enumitem}

\makeatletter
\newcommand{\applabel}[2]{\protected@edef\@currentlabel{#1}\label{#2}}
\makeatother

\title{Cosine is Nearly Human: The Reproducibility Ceiling of Perceptual Similarity}

\author{%
  Benjamin TenWolde \\
  Hal{\i}c{\i}o\u{g}lu Data Science Institute\\
  University of California, San Diego\\
  \texttt{betenwolde@ucsd.edu} \\
  \And
  Virginia R. de Sa \\
  HDSI and Cognitive Science\\
  University of California, San Diego\\
  \texttt{desa@ucsd.edu}
}

\begin{document}

\maketitle

\begin{abstract}
The NIGHTS dataset is widely used to test whether image-similarity metrics agree with human judgments of mid-level visual similarity. DreamSim, the model released with NIGHTS, reported large gains over DINO-based baselines, and those gains have been read as evidence that training on NIGHTS makes visual representations more human-aligned. We argue that the benchmark, as released, does not support that reading. Its evaluation set keeps only the triplets on which every voter agreed, about a fifth of a 100{,}000-triplet collection, and drops the many triplets where people split. We re-evaluate twelve perceptual-similarity models on the full released corpus. On its 64{,}420 non-tie triplets, cosine distance between frozen DINO features reaches 74.1\% accuracy with no training at all. Nearly half of these triplets were judged by only three or four people, which limits what any accuracy number can show. If each triplet were voted on again by a new panel of the same size, the original majority would come back only about 89\% of the time overall, and only 74--84\% of the time on the split triplets the public release leaves out. Gains from training on NIGHTS therefore show that a model fits the votes that were collected, not that it would better match a new group of people; the problem lies in how the benchmark was built, not in ordinary overfitting. The two human tasks in NIGHTS measure different things: 2AFC asks which candidate is more similar to the reference, while JND asks whether each candidate is the same as or different from it. Among triplets with both labels, JND voters judged both candidates different in 77\%, even though 2AFC voters unanimously preferred one of them. Fine-tuning also makes models more fragile under rotation, color, blur, and foreground/background changes; since no human judgments exist for these altered images, we report this as a property of the models, not a gap from humans. Plain cosine is not a universal answer either: on BAPPS, LPIPS beats every zero-shot DINO backbone on color, deblurring, and super-resolution distortions. NIGHTS supports a narrower claim than it has been used for: cosine distance on self-supervised features captures much of the reproducible signal in its current labels. Stronger claims about human alignment need a new benchmark with more votes per triplet, questions about one visual attribute at a time, per-rater data, and robustness tests.
\end{abstract}

\section{Introduction}

We make the claim that plain cosine distance on a self-supervised DINO-based backbone captures most of the human perceptual-similarity signal currently measurable in the image-triplet judgments from NIGHTS~\citep{fu2023dreamsim}. The competing claim is that supervised triplet training on NIGHTS (DreamSim and successors) yields strictly more human-aligned representations: DreamSim, a LoRA-fine-tuned DINO+CLIP+OpenCLIP ensemble trained on NIGHTS, reports 96.2\% on the public 20{,}019-triplet release and subsequent work targets that number. We argue that the cosine baseline gives the narrower interpretation because NIGHTS, the benchmark that should adjudicate whether supervised triplet training improves human alignment, is itself filtered. The 20{,}019-triplet release is the unanimous-agreement subset of a 100{,}000-triplet collection; 35{,}580 ties (50/50 splits after aggregation) and roughly 44{,}000 weak-majority triplets (patterns such as 6-1, 5-1, 3-1, 2-1) are discarded. We note that human annotators disagreed on many triplets, and no further analysis or evaluation was published on these disagreements. Annotator-level data is also missing, and the benchmark's binary labels do not distinguish contested triplets from unanimous ones. We argue that the cosine baseline (such as provided with base DINO) is therefore not just a strong starting point, but it supplies most of the signal these labels can express. No stronger alignment claim trained against the same labels can be certified on them. The remaining gap is a labeling problem, and the successor benchmark required to move past it is what this paper motivates.

This is, to our knowledge, the first evaluation on the full NIGHTS corpus. We evaluate the 64{,}420 non-tie triplets by 5-fold cross-validation, training on 51{,}536 and testing on 12{,}884 per fold. We exclude all 35{,}580 ties from training and hold them out to audit bias: 35{,}569 are 1-1 and 11 are 2-2, since voting stopped at the first disagreement. NIGHTS provides only vote tallies, so we treat votes as anonymous ballots and resample each triplet's own votes to ask how often a fresh panel of the same size would reproduce its majority. This gives a reproducibility ceiling that accounts for the noise in the small per-triplet vote counts (47\% of non-tie triplets have only 3 or 4 voters). We stratify test accuracy by vote pattern (X-0, X-1, X-2) and check left- versus right-candidate preference on tied votes.

Three findings support this claim.  (i) Resampling the per-triplet votes gives a closed-form reproducibility ceiling for any 2AFC benchmark: on NIGHTS a fresh panel of the same size reproduces the observed majority about 89\% of the time overall, and only 74--84\% on the 2-1 and 3-1 triplets that make up 46\% of the corpus. (ii) Cosine on the CLS embedding of a frozen DINOv1 backbone reaches 74\% with no training. LoRA on NIGHTS-20k raises clean accuracy by 4--6 percentage points (pp) but increases foreground-only fragility from 2.0 to 7.7 pp. The tuned models rely on background cues humans do not need to recognize objects~\citep{xiao2020backgrounds}. (iii) The two NIGHTS tasks ask different questions: 2AFC asks which candidate is more similar, and JND asks whether each candidate is the same as the reference. Among the 608 shared triplets, JND voters judged the rejected candidate different in 94\%, the chosen candidate different in 82\%, and both different in 77\%. A unanimous 2AFC choice can therefore favor one candidate even when JND voters find both different from the reference. These results show that cognitive heads (Tversky, GCM, TNN) and fine-tuning do not beat plain cosine by more than the label noise can certify. Plain cosine on DINO-family backbones is the strongest zero-shot model. The appendix lists full per-fold numbers, additional ablations (harmonization, C-RADIOv4 distillation, Qwen3-VL-as-judge), BAPPS cross-benchmark sanity checks, paired significance tests, ImageNet/non-ImageNet splits, and qualitative triplet traces (Appendix~\ref{app:traces}).

\section{Related work}

LPIPS~\citep{zhang2018lpips}, generalizing pixel-domain SSIM~\citep{wang2004ssim}, established deep-feature cosine distance as a perceptual metric. Subsequent work extended this with adversarially certified~\citep{lipsim2024} and adversarially trained~\citep{croce2025robustclip} variants. DreamSim~\citep{fu2023dreamsim} introduced synthetic Stable Diffusion v1.4 (SD)~\citep{rombach2022stablediffusion} triplets and LoRA fine-tuning~\citep{hu2021lora} for human alignment, reporting 96.2\% on the filtered 20k and becoming the default backbone for follow-up work including DiffSim~\citep{song2025diffsim} and the perceptual-alignment transfer analysis of \citet{sundaram2024perceptual}. All of these evaluations use the same 20k filtered split as the original DreamSim paper.

Self-supervised backbones are the strongest perceptual base. \citet{caron2021dino} showed that DINO's self-attention recovers object structure without supervision. \citet{muttenthaler2022vice}, \citet{muttenthaler2022human}, and \citet{conwell2024large} found that SSL ViTs align with human similarity and brain responses above supervised CNNs, complementing earlier evidence that ImageNet~\citep{deng2009imagenet} accuracy alone is a poor predictor of perceptual similarity~\citep{kumar2022perceptual}. \citet{oquab2023dinov2} scaled this to DINOv2 and DINOv3, and \citet{chen2025midlevel} benchmarked 22 SSL backbones on mid-level vision tasks including NIGHTS. A second line of work injects human alignment at training time: AligNet~\citep{muttenthaler2025alignet} learns an affine transform into a THINGS-derived human similarity space and distills it into vision backbones, and UniSim~\citep{ghazanfari2025unisim} trains a multi-task metric over 25 datasets including NIGHTS. A third line argues that attribute-similarity metrics (LPIPS, CLIP~\citep{radford2021clip}, DINO, DreamSim) miss \emph{relational} similarity, the shared logic between visually different scenes, and proposes a VLM-based metric for this axis~\citep{le2025relsim}. Across all of these approaches the final similarity operator is cosine distance (or a dot product on unit-norm embeddings), as is the case in retrieval and recommendation systems where cosine dominates despite its simplicity~\citep{steck2024cosine,you2025semantics}. We do not propose a new alignment method. We audit what NIGHTS can and cannot tell us about any model's human alignment, and we show that cosine distance captures most of the reliable agreement in NIGHTS labels.

Cognitive theories of similarity provide alternative inductive biases. Tversky's~\citep{tversky1977} contrast model uses asymmetric feature overlap in place of distance. Nosofsky's generalized context model adds selective dimensional attention~\citep{nosofsky1986}. Gentner's structure-mapping theory aligns relational representations~\citep{gentner1983}. \citet{doumbouya2025tversky} recently introduced a fully differentiable Tversky Neural Network. We evaluated these cognitive architectures directly against plain cosine-distance heads.

\section{Method}

\subsection{Dataset reconstruction}

The public NIGHTS release contains 20{,}019 triplets with binary vote columns. This is the unanimous subset of a larger 100{,}000-triplet collection that the authors released with raw vote proportions and per-triplet vote counts. The authors confirm that participant-level data was not retained, so we could not analyze individual variability across the 100{,}000 triplets. We converted the proportions to integer counts, computed the agreement pattern (larger-smaller, e.g., ``6-1'' for 6 left and 1 right), and partitioned the corpus into 64{,}420 non-tie triplets for training and evaluation and 35{,}580 tie triplets held out for the bias audit. DreamSim's trained and published results are on the  20k subset. This is the union of the X-0 (unanimous) patterns in the non-tie set. A few outlying vote counts mean that analyses restricting the unanimous subset report $N$ between 20{,}005 and 20{,}018. The vote-count distribution is skewed to small counts: 28.9\% of non-tie triplets have 3 votes, 18.2\% have 4, and 5.5\% have at least 9.

\subsection{Evaluation protocol}

We use five-fold cross-validation with the commonly used seed 42 on the 64,420 non-tie triplets. Each fold has 12,884 test triplets and 51,536 training triplets. For LoRA, we hold out 10\% of training for validation. We report mean and standard deviation of test accuracy across five folds by agreement pattern, and predictions on each fold's disjoint 7,116-triplet share of the 35,580 held-out ties.


\paragraph{Reproducibility ceiling.} We estimate how often an independent panel of voters would reproduce a triplet's majority label by resampling its observed votes, assuming exchangeable voters (\S\ref{sec:limitations}, vi). For each non-tie triplet with $L$ left votes and $R$ right votes ($V = L+R$), we draw $V$ votes with replacement from the observed proportions and ask whether the new panel's majority matches the observed one; a tie in the new panel counts half. The binomial form gives this probability in closed form per $(L,R)$ pattern. Averaging over 64{,}420 triplets gives the overall ceiling. This is an upper bound on accuracy these labels can certify: agreement beyond it may fit votes that would not recur. It does not bound accuracy against the fixed observed labels, which is trivially 100\%.

\subsection{Models evaluated}

A model is a backbone, an adaptation (frozen, LoRA, retrained, or released weights), and a head that turns embeddings into a similarity; names follow the template \texttt{\{Backbone\}-\{Adaptation\}-\{Head\}-\{Dataset\}}. Table~\ref{tab:naming} (appendix) lists the factor values. Accuracies are on NIGHTS-100k unless stated otherwise. For zero-shot models we used cosine distance on the CLS token. For LoRA models we trained rank-8 adapters on every transformer block together with the head (full recipe in Appendix~\ref{app:training}). The four heads share the same LoRA budget: plain cosine distance, our Cog head (cosine alignment, GCM context, and Tversky contrast), the inverse-frequency-stratified Cog, and the Tversky Neural Network of \citet{doumbouya2025tversky}. DreamSim-Pretrained used the authors' public weights; DreamSim-Scratch retrained the full ensemble on each fold's training set with the published codebase unchanged, giving one checkpoint per fold.

\section{Results}

\subsection{Main accuracies and the cost of unfiltering}

Table~\ref{tab:main} reports 5-fold mean accuracy on the full NIGHTS-100k for all primary models. Per-fold numbers are in Table~\ref{tab:perfold} of the appendix. Every model lost between 10.6 pp (LPIPS) and 17.2 pp (DreamSim-Pretrained) when evaluated on the full corpus instead of the unanimous X-0 subset comparable to the published 20k: on unanimous triplets the strongest models were near ceiling, while on weak-majority triplets they were near chance (Figure~\ref{fig:drop}, appendix). NIGHTS-20k kept only the unanimous triplets, so the published scores were measured on the easy part of the data.

\begin{table}[ht]
\caption{5-fold CV accuracy for selecting the human-majority candidate on NIGHTS-100k non-tie ($N=12{,}884$ per fold). $^\dagger$DreamSim-Pretrained uses released weights trained on the NIGHTS-20k training split; some of those triplets occur in our test folds, inflating its score. DreamSim-Scratch is retrained from scratch on NIGHTS-100k using the authors' published training recipe unchanged.}
\label{tab:main}
\centering
\small
\begin{tabular}{llcc}
\toprule
Model & Backbone + Head & Accuracy & Trainable params \\
\midrule
LPIPS-AlexNet-Pretrained & AlexNet + linear probe & $65.3 \pm 0.2$ & 0 \\
DINOv1-ZS-CLS-cos & DINOv1 ViT-B/16 & $\mathbf{74.1 \pm 0.5}$ & 0 \\
DINOv2-ZS-CLS-cos & DINOv2 ViT-B/14 & $70.6 \pm 0.4$ & 0 \\
DINOv3-ZS-CLS-cos & DINOv3 ViT-B/16 & $72.0 \pm 0.4$ & 0 \\
\midrule
DINOv1-LoRA-Cog & DINOv1 + Cog & $75.1 \pm 0.3$ & 641{,}346 \\
DINOv2-LoRA-Cog & DINOv2 + Cog & $77.1 \pm 0.5$ & 641{,}346 \\
DINOv2-LoRA-Cos & DINOv2 + Cosine & $78.4 \pm 0.4$ & 443{,}904 \\
DINOv2-LoRA-Cog-strat & DINOv2 + Cog + strat & $76.3 \pm 0.6$ & 641{,}346 \\
DINOv2-LoRA-TNN & DINOv2 + Doumbouya 2025 TNN & $77.4 \pm 0.5$ & 640{,}515 \\
DINOv3-LoRA-Cog & DINOv3 + Cog & $\mathbf{78.7 \pm 0.3}$ & 641{,}346 \\
\midrule
DreamSim-Pretrained$^\dagger$ & DINO+CLIP+OpenCLIP, 20k weights & $80.7 \pm 0.2$ & $\sim$1.8M \\
DreamSim-Scratch & Retrained on NIGHTS-100k, votes$\geq$6 & $78.7 \pm 0.7$ & $\sim$1.8M \\
\bottomrule
\end{tabular}
\end{table}

\subsection{Agreement-stratified reliability}

Figure~\ref{fig:tier_ceiling} splits test accuracy by the vote pattern of each triplet and sets it against the reproducibility ceiling of that pattern (\S\ref{sec:ceiling}). Zero-shot cosine on each DINO generation supplies most of what any model achieves at every level of agreement: DINOv1 scores 62.3\% on 2-1 triplets and 88.6\% on unanimous ones, with DINOv3 and DINOv2 one to five points behind. With no training, the three reach 84\%, 77\%, and 72\% of the above-chance accuracy of the retrained DreamSim ensemble. With rank-8 LoRA, a single DINOv3 backbone matches that three-backbone ensemble at every level (78.7\% overall for both), and DINOv2 comes within half a point. The lower-agreement triplets do not help as training data either: on fold 0, DINOv2 with LoRA scores 78.3\% when trained only on the 20{,}101 unanimous and X-1 ($X \geq 5$) triplets and 78.1\% once the other 26{,}282 are added (Table~\ref{tab:ladder}). Accuracy rises with agreement for every model, and the spread across models widens from 8.2 pp on 2-1 triplets to about 20 pp on unanimous ones (Figure~\ref{fig:agreement}, Appendix~\ref{app:perfold}).

\begin{figure}[ht]
\centering
\includegraphics[width=\linewidth]{figures/tier_ceiling_main.pdf}
\caption{Zero-shot DINO against what the labels can certify, by the vote pattern of the test triplet (five test folds pooled). Left: accuracy of zero-shot cosine on DINOv1, DINOv2, and DINOv3 (bars), of the same backbones after LoRA (circles), and of the retrained DreamSim ensemble (dashed), under the reproducibility ceiling of each pattern (black line); the hatched region above the ceiling would not be reproduced by a fresh panel. Right: share of the retrained ensemble's above-chance accuracy that each backbone reaches zero-shot (dark) and with LoRA (light). A single DINOv3 backbone with LoRA matches the three-backbone ensemble.}
\label{fig:tier_ceiling}
\end{figure}

Figure~\ref{fig:gap_ties} shows where the trained lead comes from. It is smallest where the labels are least reproducible: on 2-1 triplets a fresh panel would overturn 26\% of the labels, and the retrained DreamSim ensemble leads zero-shot DINOv1, DINOv3, and DINOv2 by 2.4, 4.2, and 5.3 pp. The lead roughly doubles on strong majorities and unanimous triplets (5.5--10.8 pp), where the labels are stable, and the unanimous slice alone supplies 37--39\% of each overall lead. The 35{,}580 tied triplets, which have no majority label, give the same reading. Every model picks the left candidate 50.1--50.8\% of the time, so none has a side bias, and DreamSim picks the same image as zero-shot DINOv1 on 78\% of ties, as it does on 2-1 triplets (69--70\% for DINOv2 and DINOv3); agreement rises to 84--89\% only on unanimous triplets. Fine-tuning changes behavior where voters already agreed and leaves the contested triplets nearly untouched.

\begin{figure}[ht]
\centering
\includegraphics[width=\linewidth]{figures/gap_ties_main.pdf}
\caption{Where the trained lead comes from. Left: lead of the retrained DreamSim ensemble over each zero-shot DINO backbone by vote pattern, next to the label noise of that pattern (100 minus its reproducibility ceiling). Middle: how often DreamSim-Pretrained picks the same candidate as each zero-shot backbone; tied triplets have no majority label. Right: contribution of each vote pattern to the retrained ensemble's overall lead over each backbone.}
\label{fig:gap_ties}
\end{figure}

\subsection{A reproducibility ceiling from the vote counts}
\label{sec:ceiling}

Table~\ref{tab:ceiling} reports the reproducibility ceiling for each subset: the probability that a fresh panel of the same size, drawn by resampling a triplet's observed votes with replacement, reproduces the observed majority label (a tie in the new panel counts half). The assumption is that each triplet's observed votes are exchangeable draws from one voter population (\S\ref{sec:limitations}, vi). The ceiling is a property of the labels, not of any model. We report it as a ceiling on the accuracy these labels can certify, not a bound on accuracy against the fixed observed labels: a deterministic model that memorized the observed majorities reaches 100\%. The question is what a model generalizing from those labels can be expected to reproduce on a fresh vote. The overall ceiling on NIGHTS-100k non-tie is about 89\%. The aggregate hides a split: 100\% on unanimous triplets but far less where the majority hangs on one or two votes: a fresh panel of three agrees with a 2-1 majority only 74\% of the time (20/27; $N{=}18{,}454$), a fresh panel of four with a 3-1 majority 84\% ($N{=}11{,}181$), and a fresh panel of five with a 4-1 majority 94\% ($N{=}6{,}841$). The 2-1 and 3-1 patterns together account for 46\% of the non-tie corpus and drag the ceiling down.

\begin{table}[ht]
\caption{Reproducibility ceiling on NIGHTS-100k: the probability that a fresh panel of the same size, drawn by resampling each triplet's observed votes with replacement, reproduces the observed majority label. The strong-majority row spans its patterns (5-1 through 10-1); the overall value lies between 88.8\% and 89.2\% for any mix.}
\label{tab:ceiling}
\centering
\begin{tabular}{lcc}
\toprule
Subset & N & Repeatability \\
\midrule
Unanimous (X-0) & 20{,}018 & 100.0\% \\
Strong majority (X-1, $X \geq 5$) & 7{,}922 & 96.5--99.5\% \\
Weak majority (4-1, 3-1, 2-1) & 36{,}476 & 81.0\% \\
\midrule
All non-tie (NIGHTS-100k) & 64{,}420 & $\approx$89\% \\
\bottomrule
\end{tabular}
\end{table}

No model reaches this ceiling: the best, DreamSim-Pretrained, scores 80.7\% against the $\approx$89\% the labels could certify, and on 2-1 triplets, where the label itself repeats only 74\% of the time, every model sits below that bar (57.0--65.2\%, Fig.~\ref{fig:agreement}). What the ceiling rules out is reading small gaps as better alignment: differences of a few percentage points fall inside the label noise, so they cannot be attributed to one model matching people better. The models that gained the most clean-image accuracy also lost invariance (\S\ref{sec:shape_axes}), and 96.6\% of DreamSim-Pretrained's disagreements with the majority label matched at least one dissenting voter.

The root cause is the NIGHTS collection procedure, which solicited votes until the first disagreement and then stopped. As a result, 47\% of non-tie triplets were labeled by 3 or 4 voters. For a triplet with true preference $p = 0.6$, the probability that a 3-voter sample produces a majority in the other direction is $\binom{3}{0}(0.6)^0(0.4)^3 + \binom{3}{1}(0.6)^1(0.4)^2 = 0.352$, so better than one in three observed majority labels would flip under re-annotation. A stable majority at $p=0.6$ requires on the order of 35 votes per triplet. Model differences of a few percentage points on NIGHTS-100k fall within this label-noise band. They cannot be attributed to model quality without re-annotation data.

\subsection{Cognitive priors and backbone ordering}

We asked whether a psychologically motivated similarity head beats plain cosine. We tested three, Tversky asymmetric contrast, GCM selective attention, and the differentiable TNN of \citet{doumbouya2025tversky}. None did: cosine reached 78.4\% $\pm$ 0.4 against 77.1\% for Cog, 76.3\% for inverse-frequency stratification, and 77.4\% for TNN. The Tversky head shows why. It has two free weights, $\alpha$ for features only the reference has and $\beta$ for features only the candidate has; unequal values would make similarity directional, but the fitted values came out nearly equal, so the head reduced to a symmetric measure and behaved like cosine. Retraining DreamSim from scratch on NIGHTS-100k with the authors' code gave $78.7 \pm 0.7$\%, 2.0 pp below its published weights on the same splits. Both gaps are smaller than the label noise: on 2-1 and 3-1 triplets a fresh panel would agree with the observed majority only 74--84\% of the time, and differences under $\sim$10 pp are not separable from that noise (Section~\ref{sec:ceiling}). A Cohen's-$\kappa$ inter-model analysis is in Appendix~\ref{app:kappa}.

On NIGHTS-100k, zero-shot DINOv1 (74.1\%) beat DINOv2 (70.6\%) and DINOv3 (72.0\%). Under LoRA the order reversed: DINOv3-LoRA-Cog 78.7\%, DINOv2-LoRA-Cos 78.4\%, DINOv1-LoRA-Cog 75.1\%. Zero-shot rank does not predict fine-tuned rank (DINOv3 gained 6.7 pp).

\paragraph{Backbone family and scale.} Cosine is not DINO-specific. We ran sixteen zero-shot backbones: eight CNNs (five supervised, three bio-inspired), two non-DINO ViTs (CLIP and supervised), and six DINO ViTs (DINOv1, DINOv3, and DINOv2 at S, B, L, and g). The ten non-DINO models span 65.3--72.4\%, and six of them score above 70\% (Table~\ref{tab:backbones}). Bigger is not better: DINOv2 falls monotonically over 52$\times$ in size on matched folds, and its two largest fall more than 5 pp behind DINOv1 cosine (Table~\ref{tab:scale}), extending the CLIP/MAE trend of \citet{fu2023dreamsim}. Recent large models do worse: C-RADIOv4-SO400M at 58.2\% after one fold-0 epoch and Qwen3-VL-4B at 63.0\% (Appendix~\ref{app:extra}).

\paragraph{The fine-tuning gain is real but not certifiable as alignment.}\label{sec:gain} The untrained DreamSim ensemble reached 74.5\% versus 74.1\% for frozen DINOv1 (0.4 pp from ensembling). Retraining reached 78.7\% (4.2 pp from training), while the pretrained checkpoint gained 2.0 pp overall and 3.8 pp on its overlapping 20k training subset. The 4.2 pp gain is not a color artifact: splitting by candidate color differences yields similar gains ($+3.9$ and $+4.9$ pp; Appendix~\ref{app:colorcontrol}). Training learns real signal from the labels, but their 74--84\% weak-majority reproducibility cannot establish better human agreement (\S\ref{sec:ceiling}).

\subsection{The ceiling and the ranking on a second benchmark (BAPPS)}
\label{sec:bapps}

The same resampling applies to BAPPS~\citep{zhang2018lpips}, which collects 5 votes per triplet ($N=36{,}344$). A 3-2 majority is reproduced by a fresh panel of five only 68.3\% of the time and a 4-1 majority 94.2\%, so 2AFC benchmarks with few votes per triplet carry the same per-pattern limits on what their labels can certify.

On BAPPS, the ranking reverses: LPIPS beats every zero-shot DINO backbone by 1.4 to 5.6 pp, while DINO leads on NIGHTS by 5 to 10 pp. Only DreamSim-Pretrained (BAPPS 73.6, NIGHTS 80.7) and zero-shot DINOv1 (71.5, 74.1) are competitive on both (Table~\ref{tab:bapps}). Plain cosine is not competitive for low-level distortions; no single score captures performance across both benchmarks. Per-category comparisons appear in Appendix~\ref{app:crossover}.

\begin{table}[ht]
\caption{Zero-shot 2AFC accuracy on BAPPS validation (six distortion categories) and NIGHTS-100k. LPIPS leads on BAPPS, DINOv2/v3 on NIGHTS; only DreamSim-Pretrained and DINOv1 are competitive on both.}
\label{tab:bapps}
\centering
\small
\begin{tabular}{lccc}
\toprule
Model & BAPPS & NIGHTS-100k & Above 69.0\,/\,70.1\% \\
\midrule
LPIPS-AlexNet           & 72.9 & 65.3 & \\
DINOv1-ZS               & 71.5 & 74.1 & \checkmark \\
DINOv2-ZS               & 67.3 & 70.6 & \\
DINOv3-ZS               & 68.4 & 72.0 & \\
DreamSim-Pretrained & 73.6 & 80.7 & \checkmark \\
\bottomrule
\end{tabular}
\end{table}

\subsection{Shape, texture, and the axes of perceptual alignment}
\label{sec:shape_axes}

Accuracy against NIGHTS majority labels shows how often a model agrees, not which image features it uses. We therefore replicate the DreamSim appendix's stimulus-manipulation study~\citep{fu2023dreamsim}, measuring accuracy under eleven manipulations.

\paragraph{Eleven stimulus manipulations.} We apply each manipulation identically to all three images in a triplet from the unanimous 20k subset ($N=2000$ sampled for speed), then score accuracy against the unchanged human label. The eleven manipulations cover four of the five DreamSim Figure 6 axes: orientation (90/180/270$^\circ$), color (grayscale), shading (HSV V-channel histogram equalization), structure (Gaussian blur $\sigma\in\{1,2,4,8\}$), and foreground/background (center disc of radius 0.35$\cdot$image side or its complement). NIGHTS has no human labels for manipulated images, and we collected none. These tests measure model invariance, not human perceptual agreement; per-axis results are in Appendix~\ref{app:axes_detail}.

\paragraph{The same backbone, trained longer, loses invariance.} Across 16 backbones, higher clean accuracy correlates with larger changes under manipulation, but this does not isolate training. We train DINOv2 with rank-8 LoRA on NIGHTS-20k for 8 epochs and test checkpoints \{0,1,2,3,5,8\} on the same 2000 triplets. Fragility is the mean absolute change in accuracy across the eleven manipulations, using absolute values because some increase accuracy. From epoch 0 to 1, clean accuracy rises from 84.6 to 89.8\% and fragility from 1.9 to 5.1 pp. By epoch 8, accuracy falls to 85.2\% while fragility climbs to 12.3 pp. With backbone and data fixed, training duration causes the invariance loss. The trajectory and attention overlays are in Appendix~\ref{app:extra}.

\paragraph{Foreground-only fragility rises after LoRA.} NIGHTS uses Stable Diffusion images with centered subjects and low-detail, nonuniform backgrounds. The foreground-only mask keeps a centered disc of radius 0.35 times the image side (about 38\% of its area) and sets the surround to black. We apply the same mask to all three images without cropping or rescaling, preserving subject size and position. Removing the surround tests whether the model relies on its cues despite their low detail. Zero-shot DINOv2 drops 2.0 pp; DINOv2 + LoRA drops 7.7 pp, a 3.8$\times$ increase. No human labels exist for these manipulated triplets, so this measures model invariance, not model--human divergence. In related natural-image categorization tasks, humans lose $\sim$3 pp under IN-9 background randomization~\citep{xiao2020backgrounds} and $\sim$1 pp on ObjectNet~\citep{barbu2019objectnet}. The zero-shot drop is within that range; the LoRA drop is outside it, though these are not synthetic-triplet similarity judgments.

\paragraph{Fine-tuning makes models brittle; zero-shot DINO stays robust.} LoRA-trained and DreamSim-Pretrained models are both more accurate on clean images and more fragile under the 11 manipulations than every zero-shot backbone. Each point of clean accuracy gained above the zero-shot band costs about a point of mean fragility (Appendix Figure~\ref{fig:cost_acc}), while zero-shot DINOv3 is the most invariant backbone tested. What DreamSim reports as better human alignment thus comes with more brittleness than the zero-shot ViT and DINO backbones. 

\paragraph{A 2AFC lead does not carry over to JND.} No model is trained for JND; all training uses the 2AFC labels alone. NIGHTS also ships a second human task, JND, on 608 of the 2AFC triplets: each candidate was shown briefly with the reference and judged ``same'' or ``different,'' giving 1{,}216 reference--candidate pairs labeled by majority vote. We score the same frozen backbones on it, with no retraining: we reuse the embeddings from the main results and use the JND labels only to grade them, by how well the distance a backbone puts between a pair's two images separates the ``different'' pairs from the ``same'' ones (ROC AUC, which needs no cutoff). On the 2AFC triplets DINOv1 beats DINOv2 by 3.8 pp; on JND the lead disappears (AUC 0.674 vs.\ 0.677; Figure~\ref{fig:jnd_2afc_tier}b, Table~\ref{tab:axes}). Across the fourteen zero-shot models, 2AFC accuracy barely predicts JND AUC (Spearman $\rho = 0.37$, $p = 0.19$). The two models trained on NIGHTS 2AFC score highest on both tasks (DreamSim 0.727, DINOv2+LoRA 0.720), though a single AUC carries a 95\% interval of about $\pm 0.04$. One benchmark number is thus a weak stand-in for perceptual alignment (Appendix~\ref{app:axes}).

\begin{figure}[ht]
\centering
\includegraphics[width=\linewidth]{figures/jnd_main_v4.pdf}
\caption{The two human tasks on the same images. (a) Share of candidates JND voters judged ``different'' from the reference, by the triplet's 2AFC vote count: 94\% of rejected and 82\% of chosen candidates overall. (b) 2AFC accuracy against JND AUC for sixteen backbones: among zero-shot models (filled) the two are only weakly related, and DINOv1's 2AFC lead disappears. $^\ddagger$This checkpoint's training fold contained 80\% of the JND triplets; on the 244 pairs it never saw, its AUC is 0.684, inside the zero-shot range. (c) One triplet per case: both candidates judged different (77\% of triplets), only the chosen one judged the same (103 triplets), only the rejected one judged the same (34). Full breakdown in Appendix~\ref{app:disagree}.}
\label{fig:jnd_2afc_tier}
\end{figure}

\subsection{Model confidence versus the human vote split}
\label{sec:confidence}

Accuracy only shows whether a model picks the same image as most people. It does not ask whether the model hedges when people are also split on the choice.  To check that, we turn each model's score into a probability of choosing the left image and read how confident it is (the calibration and tables are in Appendix~\ref{app:confidence}). No model is trained to output confidence; it is read off afterward with the same calibration for every model, so the differences come from the models, not the calibration. The human vote sets the target: a tie means no preference, a unanimous triplet full agreement.

The baseline is untrained cosine, and fine-tuning makes confidence worse, not better. On the 35{,}580 tied triplets, where the vote supports zero confidence, untrained DINOv1 cosine already reports 0.38 (LPIPS 0.25); fine-tuning raises this to 0.48 for DreamSim-Pretrained and to 0.80--0.95 for our own LoRA heads. Even on 2-1 triplets DreamSim-Pretrained sits at 0.49, and on weak-majority triplets, where the majority hangs on one or two votes, 11.8\% of its predictions still exceed 0.9. On unanimous triplets the same upward shift helps, since the target is high (DreamSim 0.79 against untrained DINOv1's 0.58). So the overconfidence grows with training, measured against the untrained cosine it starts from. 

Fine-tuning does have one edge, and we do not dispute it. A trained model's confidence rises and falls with the human vote more tightly than an untrained one: the two move together with rank correlation 0.76 for DreamSim-Pretrained, against 0.62 for zero-shot DINOv1 and 0.41 for LPIPS. So training teaches a model to follow how much people agree, without teaching it to be the right amount of unsure when they are split. Tables for each agreement level and the reliability curves are in Appendix~\ref{app:confidence}.

\section{Discussion}
\label{sec:discussion}

\paragraph{Why cosine is nearly human, and why more training is not quite there.} Plain cosine distance on the CLS embedding of a frozen DINO backbone reached 74\% on NIGHTS-100k with no training, and its accuracy drop under the foreground-only mask stayed inside the range external human anchors suggest. LoRA fine-tuning raised clean accuracy from 84.6 to 89.3\%, but that 4.7-pp gain came with a loss of invariance: the foreground-only drop grew from 2.0 to 7.7 pp while the object was fully visible, and DreamSim-Pretrained was the most accurate and the most fragile model across orientation, color, heavy blur, and foreground and background masks (Fig.~\ref{fig:fig6_replica}). With the same LoRA budget, no cognitive head (Tversky, GCM, TNN) beat plain cosine, and the Tversky heads learned no direction, so they reduced to symmetric measures. The 4.2-pp gain from training survives every control we ran. But NIGHTS cannot tell us whether the gain reflects better agreement with people in general or a closer fit to how these voters weighted features in these images. The labels are too noisy to decide: a fresh panel reproduces the majority only 81\% of the time on weak-majority triplets and 74\% on 2-1, so a 5-pp gap between models falls inside that noise. The JND conflict makes the same point without any resampling (Fig.~\ref{fig:jnd_2afc_tier}). The 96.2\% DreamSim reports is measured on the 20\% of triplets where voters already agreed unanimously, where the ceiling is 100\% by construction. Retraining the same ensemble on the full 100k drops it to 94.1\% on that unanimous subset, 3.8 pp below the 97.9\% it reaches when trained on the subset alone: the published weights are tuned to exactly the easy slice the 20k release keeps. We therefore read DreamSim not as a model that matches people better, but as one fit to the specific voters and images in the 20k set. To distinguish broader alignment from this kind of fit, we need labels that record which features each voter used and models that compare images feature by feature rather than reducing them to one distance.

\paragraph{NIGHTS is a single-axis benchmark presented as a general one.} Perceptual similarity spans three levels: low-level pixel fidelity (where LPIPS and BAPPS operate), mid-level attributes of an object (pose, color, shape, texture, lighting), and high-level object identity and scene layout. NIGHTS sits between the mid and high levels: a voter combines several attributes and an overall judgment into one binary choice. We hypothesize that unanimous triplets differ mainly on one mid-level feature, while trade-offs across features drive the weak-majority votes in the 80{,}000 triplets NIGHTS-20k excluded. The 16-pp gap between DreamSim's 96.2\% on 20k and 80.7\% on 100k is the cost of leaving that single-axis subset. The reproducibility ceiling reflects more than random noise. Voters may weigh features of the same images differently, but the 20k release hides these disagreements by keeping only unanimous triplets. A model tuned to one benchmark inherits that benchmark's particular set of voters, so scoring well on it is not the same as matching people in general. 

\paragraph{Voting is preference, not an explanation.} Both NIGHTS 2AFC (``which is more similar?'') and NIGHTS-JND (``same or different?'') collect a single bit per voter. Multi-axis differences (shape, color, pose, lighting, composition, background, identity) collapse into one vote that records the choice but not the attribute the voter used. The two labels can be compared directly. NIGHTS-JND reuses 608 unanimous 2AFC triplets, so each of their 1{,}216 reference--candidate pairs carries a 2AFC label (its candidate won or lost) and a JND label (``same'' or ``different'' from the reference). JND majorities judged 94\% of rejected candidates and 82\% of chosen candidates different from the reference, a pattern consistent with the 2AFC preference (Figure~\ref{fig:jnd_2afc_tier}). The difference varies with the number of 2AFC votes; both rates round to 91\% in the 10--0 group, which contains only 22 triplets. In 77\% of triplets, JND majorities judged both candidates different. Among the 137 triplets in which exactly one candidate was judged ``same,'' it was the rejected candidate in 34 (25\%), and the chosen candidate in 103 (75\%). Both tasks also split internally: 47\% of JND pairs and 69\% of non-tie 2AFC triplets have a dissenting voter (Appendix~\ref{app:disagree}). A 2AFC vote is a forced comparison with no ``neither'' option, so a candidate can win unanimously by being the less different of two clearly different images. A unanimous 2AFC label thus means ``closer than the alternative,'' not ``similar to the reference.'' Whether a shared background should count toward similarity is itself one of these collapsed axes: the fragility results in \S\ref{sec:shape_axes} show that tuned models weigh it more heavily than untuned ones, but without per-rater attribute labels we cannot tell whether human voters weighed it at all. The tasks use the same image pairs but ask different questions, so the labels capture different aspects of similarity. Without per-voter labels, we cannot tell whether people perceived the images differently or interpreted the questions differently. The same release omits annotator identifiers and demographics.

\subsection{A protocol for the next benchmark}
\label{sec:protocol}

Models are not the bottleneck; the labels are. Today's vision backbones already do the things perceptual similarity needs: they match points on an object as its pose changes~\citep{amir2021deepvit,caron2021dino}, stay stable when the view rotates~\citep{naseer2021intriguing,raghu2021vits}, can be tuned to judge by shape rather than texture~\citep{geirhos2018imagenettrained,hermann2020origins}, and their sensitivity to color, rotation, and blur can each be measured (Appendix~\ref{app:shape_axes_extra}). What they cannot get past is a label that records only one forced choice per triplet and discards which feature the person used to make it. A benchmark that fixed this would let us ask whether a model is right for the right reason. It would: \textbf{(1)} change one feature at a time, by taking a single image and editing it by a small and a large amount along one dimension such as brightness or blur, so the correct answer is fixed without asking anyone; \textbf{(2)} ask the same question NIGHTS asked (``which image is more like the center one?'') without naming the dimension, and after the choice have the person mark on the image what looked different, say where, and rate how sure they were, so each vote comes with its reason; \textbf{(3)} release every person's own responses, not just the vote totals, so a model's mistakes can be checked against how people actually split; \textbf{(4)} collect enough votes per triplet to push the reproducibility ceiling of \S\ref{sec:ceiling} close to 100\%, so a reported gap between models means something; \textbf{(5)} fix the train/test splits and tie rules and ship them with the data, so every model is scored the same way; and \textbf{(6)} report accuracy on held-out image distortions next to the headline number, so robustness is visible.

A first study re-annotates 3{,}000 NIGHTS triplets stratified by agreement pattern, with 30 annotators per triplet, no early stopping, and per-rater votes in the release. A second condition has the same annotators judge the \S\ref{sec:shape_axes} manipulations of those triplets, supplying the human baseline the invariance results currently lack, and the 608 triplets shared with NIGHTS-JND, so the 2AFC--JND dissociation can be traced to raters. Simply re-annotating the same images would not help: more votes would only shrink the error bars, not answer the finer questions a multi-axis test is built to ask. 

\section{Limitations}
\label{sec:limitations}

\textbf{Scope caveats.} \textbf{(i)} The model accuracies in Table~\ref{tab:main} are specific to NIGHTS, although the reproducibility ceiling itself transfers to other 2AFC benchmarks (\S\ref{sec:bapps}). \textbf{(ii)} We ran sixteen backbones zero-shot but fine-tuned only the three DINO models, and the size sweep covers one family (DINOv2, S to g); other families may scale differently. \textbf{(iii)} All LoRA models used rank 8 so that the heads were compared at equal adapter capacity; we did not sweep the rank. \textbf{(iv)} When DreamSim-Pretrained disagreed with the majority label, a minority voter had made the same choice in 96.6\% of cases; without per-voter labels we cannot tell whether those triplets were ambiguous or whether the model shares a bias with those voters. \textbf{(v)} The claim that some JND/2AFC disagreements reflect different readings of the instructions rests on one co-author looking at ten JND pairs, not a systematic re-annotation, so we treat it only as a hypothesis. \textbf{(vi)} The reproducibility ceiling resamples each triplet's observed votes assuming exchangeability, taking the observed vote fraction at face value, which is optimistic for small panels; different pools, instructions, and within-session effects would lower real agreement. Independent re-annotation reproducing the labels well above it would overturn our reading. \textbf{(vii)} Public NIGHTS votes lack rater IDs, demographics, per-rater consistency, and quality filters, so rater bias, subgroup accuracy, and majorities within a representative subpopulation cannot be tested~\citep{henrich2010weirdest}.

\section*{Broader impact and reproducibility}

We also release the seed-42 5-fold definitions, as well as  per-triplet test predictions for all models.  We release the retrained checkpoints, the standalone reproducibility-ceiling script parameterized by vote count, and the analysis notebooks that regenerate every table and figure from the released predictions, at \url{https://github.com/benlten/cosine_is_nearly_human} under the MIT license. All human judgment data were obtained from the public NIGHTS, NIGHTS-JND, and BAPPS releases.  There was no collection of new human data.

\section*{Acknowledgements and disclosures}

\textbf{Funding.} Benjamin TenWolde is supported by the Department of Defense (DoD) SMART Scholarship-for-Service Program. This work used hardware originally provided by funding from Sony and NSF IIS-CRCNS-2208362.

\textbf{Competing interests.} The authors declare no competing interests.

\bibliographystyle{plainnat}
\bibliography{references}

% !TEX root = cosine_is_human.tex
\appendix
\clearpage

\noindent{\Large\bfseries A\quad Model naming and factor values}
\applabel{A}{app:naming}
\vspace{1em}

Throughout the paper we identify each evaluated model by a four-part name of the form \texttt{\{Backbone\}-\{Adaptation\}-\{Head\}-\{Dataset\}}. Table~\ref{tab:naming} enumerates the possible values each component can take. Backbones are the feature extractors. Adaptation describes whether the backbone is frozen, partially trained via LoRA adapters, fully retrained, or used with externally-supplied pretrained weights. The head is the module that converts two (or three) backbone outputs into a similarity score; nine variants are compared including plain cosine distance, three cognitive-science-motivated architectures (Cog-v1 through Cog), the published Tversky Neural Network of \citet{doumbouya2025tversky}, the DreamSim margin-hinge ensemble loss, and a VLM logit-ratio judge. Dataset values record which NIGHTS subset a model was trained and evaluated on; omitting the dataset suffix defaults to NIGHTS-100k, our primary evaluation target.

\begin{table}[!ht]
\caption{Model naming scheme. Every model name in the paper is a concatenation of one value from each of the four rows. Example: \texttt{DINOv3-LoRA-Cog} uses DINOv3 as backbone, rank-8 LoRA on qkv/proj, our Cog head (frozen cosine alignment + GCM context + Tversky contrast), evaluated on the full NIGHTS-100k corpus.}
\label{tab:naming}
\centering
\small
\begin{tabular}{lll}
\toprule
Component & Values & Description \\
\midrule
Backbone & DINOv1, DINOv2, DINOv3 & Self-supervised ViTs \\
 & C-RADIOv4 & Distillation of DINOv3+SigLIP2+SAM3 \\
 & Qwen3-VL-4B & Vision-language model \\
 & DreamSim & DINO+CLIP+OpenCLIP ensemble \\
 & AlexNet-lin & LPIPS classical backbone \\
\midrule
Adaptation & ZS & Zero-shot, frozen backbone \\
 & LoRA & Rank-8 LoRA on qkv/proj + head trained \\
 & Pretrained & External weights (trained on NIGHTS-20k) \\
 & Scratch & Retrained on the stated dataset from scratch \\
\midrule
Head & CLS-cos & Cosine distance on CLS token \\
 & Tversky-frozen & Tversky contrast, $\alpha=\beta=0.5$ \\
 & Cog-v1 & Learned cross-attention + context + Tversky \\
 & Cog-v2 & Shared alignment + bottleneck + Tversky \\
 & Cog & Frozen cosine align + GCM context + Tversky \\
 & Cog-strat & Cog with inverse-frequency sampling \\
 & TNN & Doumbouya 2025 differentiable TNN \\
 & Hinge-ens & Margin hinge loss over ensemble features \\
 & VLM-judge & Qwen3-VL logit ratio on A/B choice \\
\midrule
Dataset & NIGHTS-100k & Full 100k, 64{,}420 non-tie, 35{,}580 ties held out \\
 & NIGHTS-100k-unan & X-0 subset, $N{=}20{,}017$ \\
 & NIGHTS-100k-weak & 2-1/3-1/4-1 subset, $N{=}36{,}476$ \\
 & NIGHTS-20k (published) & Filtered unanimous release \\
\bottomrule
\end{tabular}
\end{table}

\clearpage
\noindent{\Large\bfseries B\quad Per-fold raw numbers}
\applabel{B}{app:perfold}
\vspace{1em}

Table~\ref{tab:main} reports mean accuracy and standard deviation over five cross-validation folds. Table~\ref{tab:perfold} gives the accuracy of each of the twelve models on each fold. We form folds from the 64{,}420 non-tie triplets using fixed seed 42 and reuse the same indices for every model, so comparisons are not affected by data shuffling. For trained models, no fold differs from its model's mean by more than 1 percentage point, and the model rankings stay the same on every fold. Figure~\ref{fig:drop} plots each model's accuracy on the unanimous X-0 subset against its accuracy on the full corpus: every model drops 10.6--17.2 pp as the discarded disagreement is reincorporated.

\begin{figure}[ht]
\centering
\includegraphics[width=0.78\linewidth]{figures/drop_chart.pdf}
\caption{The same models use 5-fold cross-validation on unanimous X-0 test folds (blue; disjoint from training, comparable to published 20k) and all non-tie NIGHTS-100k triplets (red). Each drops 10 to 17 points because test sets differ. Dashed line: DreamSim's published 96.2\% on their 20k held-out split.}
\label{fig:drop}
\end{figure}

Figure~\ref{fig:agreement} gives the per-model curves behind Figures~\ref{fig:tier_ceiling} and~\ref{fig:gap_ties}, and Table~\ref{tab:tiergap} the accuracies by vote pattern.

\begin{figure}[!ht]
\centering
\includegraphics[width=0.78\linewidth]{figures/agreement_curve.pdf}
\caption{Accuracy versus human-agreement strength on NIGHTS-100k non-tie triplets. Accuracy rises with agreement for every model; DreamSim leads at most levels, but its released weights overlap our test folds. Unanimous patterns are pooled at 1.0; patterns with fewer than 100 triplets are omitted. Shaded band: the unanimous NIGHTS-20k subset.}
\label{fig:agreement}
\end{figure}

\begin{table}[!ht]
\caption{Accuracy by agreement tier, pooled over the five test folds, against the per-tier reproducibility ceiling (\S\ref{sec:ceiling}); the numbers behind Figures~\ref{fig:tier_ceiling} and~\ref{fig:gap_ties}. The LoRA columns are DINOv1-LoRA-Cog, DINOv2-LoRA-Cos, and DINOv3-LoRA-Cog. $^\dagger$Released weights overlap our test folds (Table~\ref{tab:main}).}
\label{tab:tiergap}
\centering
\footnotesize
\setlength{\tabcolsep}{4pt}
\begin{tabular}{lcc ccc ccc cc}
\toprule
 & & & \multicolumn{3}{c}{Zero-shot cosine} & \multicolumn{3}{c}{+ LoRA} & \multicolumn{2}{c}{DreamSim} \\
\cmidrule(lr){4-6}\cmidrule(lr){7-9}\cmidrule(lr){10-11}
Agreement & N & Ceiling & v1 & v2 & v3 & v1 & v2 & v3 & Pre$^\dagger$ & Scratch \\
\midrule
2-1 & 18{,}454 & 74\% & 62.3 & 59.4 & 60.5 & 61.7 & 63.9 & 64.1 & 65.2 & 64.7 \\
3-1 & 11{,}181 & 84\% & 67.6 & 64.9 & 66.5 & 68.3 & 72.1 & 72.0 & 73.0 & 71.9 \\
4-1 & 6{,}841 & 94\% & 71.5 & 68.3 & 69.8 & 72.8 & 77.3 & 77.7 & 79.5 & 77.2 \\
X-1 ($X \geq 5$) & 7{,}922 & 96.5--99.5\% & 76.2 & 72.7 & 73.8 & 78.8 & 82.7 & 83.0 & 85.2 & 83.5 \\
X-0 (unanimous) & 20{,}019 & 100\% & 88.6 & 83.9 & 85.7 & 90.7 & 93.8 & 94.5 & 97.9 & 94.1 \\
\midrule
All non-tie & 64{,}420 & $\approx$89\% & 74.1 & 70.6 & 72.0 & 75.1 & 78.3 & 78.7 & 80.7 & 78.7 \\
\bottomrule
\end{tabular}
\end{table}

\paragraph{Training-tier ladder.} Table~\ref{tab:ladder} asks whether the lower-agreement triplets help as training data. We trained DINOv2-LoRA-Cos on nested subsets of the fold-0 training split, ordered by agreement, and scored every run on the same fold-0 test split. A common validation set (10\% of the fold-0 training split) was removed from every run, and each run kept the epoch with the best validation accuracy on the tiers it trained on. Training on unanimous triplets alone gives 94\% of the gain over zero-shot that training on all triplets gives. Adding the remaining 21{,}380 triplets, almost all 2-1 and 3-1, to the first three tiers does not raise accuracy (78.4\% to 78.1\%). At equal size, unanimous triplets beat a random sample by 0.8 pp. In every run, validation accuracy peaked after one epoch, and by epoch 10 training accuracy had reached 100\% and test accuracy had fallen below zero-shot. These are single runs on one fold: an accuracy in the table has a standard error of 0.4 pp overall and 0.7--1.2 pp within a tier.

\begin{table}[!ht]
\caption{Training-tier ladder on fold 0. DINOv2-LoRA-Cos is trained on nested subsets of the training split and scored on the fold-0 test split (12{,}884 triplets) by vote pattern. The first row is the same backbone before training; the last is a random sample of all non-tie triplets at the size of the unanimous set.}
\label{tab:ladder}
\centering
\small
\begin{tabular}{lrcccccc}
\toprule
Training set & Triplets & 2-1 & 3-1 & 4-1 & X-1 ($X \geq 5$) & X-0 & All \\
\midrule
None (zero-shot) & 0 & 59.9 & 64.5 & 68.2 & 72.4 & 84.5 & 70.8 \\
X-0 & 14{,}459 & 63.7 & 71.4 & 76.2 & 81.0 & 93.1 & 77.7 \\
X-0, X-1 ($X \geq 5$) & 20{,}101 & 64.8 & 72.2 & 77.3 & 80.8 & 93.4 & 78.3 \\
X-0, X-1 ($X \geq 5$), 4-1 & 25{,}003 & 64.5 & 71.4 & 77.6 & 82.2 & 93.7 & 78.4 \\
All non-tie & 46{,}383 & 62.9 & 71.9 & 78.9 & 82.4 & 93.5 & 78.1 \\
\midrule
Random, X-0 size & 14{,}459 & 63.4 & 71.1 & 74.6 & 81.2 & 91.6 & 76.9 \\
\bottomrule
\end{tabular}
\end{table}


\begin{table}[!ht]
\caption{Per-fold accuracies (\%) on the NIGHTS-100k non-tie corpus ($N{=}64{,}420$ triplets; 12{,}884 per test fold; 5-fold cross-validation, seed 42). Each row is one model identified by the naming scheme of Table~\ref{tab:naming}. Columns are the five test folds; the same fold definitions are reused across all models. Mean and standard deviation over the five folds are reported in Table~\ref{tab:main} of the main body.}
\label{tab:perfold}
\centering
\small
\begin{tabular}{lccccc}
\toprule
Model & Fold 1 & Fold 2 & Fold 3 & Fold 4 & Fold 5 \\
\midrule
LPIPS-AlexNet-Pretrained & 65.4 & 65.5 & 65.3 & 65.3 & 65.0 \\
DINOv1-ZS-CLS-cos & 74.8 & 74.0 & 73.8 & 73.6 & 74.2 \\
DINOv2-ZS-CLS-cos & 70.8 & 71.0 & 70.8 & 70.3 & 70.0 \\
DINOv3-ZS-CLS-cos & 72.4 & 72.3 & 72.1 & 71.4 & 71.7 \\
DINOv1-LoRA-Cog & 74.8 & 75.6 & 75.0 & 75.3 & 74.9 \\
DINOv2-LoRA-Cog & 77.4 & 77.5 & 76.2 & 77.4 & 77.1 \\
DINOv2-LoRA-Cos & 78.7 & 78.4 & 77.6 & 78.3 & 78.8 \\
DINOv2-LoRA-Cog-strat & 76.4 & 76.4 & 76.2 & 75.8 & 77.2 \\
DINOv2-LoRA-TNN & 77.2 & 77.5 & 76.7 & 77.4 & 78.1 \\
DINOv3-LoRA-Cog & 78.7 & 78.9 & 78.3 & 78.4 & 79.1 \\
DreamSim-Pretrained & 80.5 & 80.9 & 80.4 & 80.9 & 80.8 \\
DreamSim-Scratch & 79.4 & 79.7 & 77.9 & 78.4 & 78.2 \\
\bottomrule
\end{tabular}
\end{table}

\clearpage
\noindent{\Large\bfseries C\quad Full Cohen's $\kappa$ matrix}
\applabel{C}{app:kappa}
\vspace{1em}

Cohen's $\kappa$ quantifies the agreement between two classifiers after correcting for the agreement expected by chance. We compute $\kappa$ pairwise between eleven of the twelve models in Table~\ref{tab:main} (the TNN head is omitted) on the pooled 64{,}420 non-tie test predictions (every triplet appears in exactly one fold's test split, so the pooled set is a disjoint union of the five held-out folds). Values near 1 indicate redundant models that would gain little from ensembling; values near 0.4--0.6 indicate models that make different errors and are candidates for ensembling. In Figure~\ref{fig:kappa} the three DINO zero-shot entries cluster tightly ($\kappa>0.63$), as do the five LoRA-adapted DINO variants ($\kappa>0.6$), confirming that the DINO family makes highly correlated errors. LPIPS sits further out ($\kappa=0.25\text{--}0.36$ with any DINO-family model), which is consistent with its different underlying feature basis. DreamSim-Pretrained and DreamSim-Scratch reach $\kappa=0.75$, the highest pairwise agreement in the matrix: retraining the DreamSim ensemble on NIGHTS-100k does not shift its error structure, it only shifts the overall accuracy by 2 pp. The scratch model is structurally the same predictor, trained on a cleaner subset, that fails on the same ambiguous triplets.

\begin{figure}[!ht]
\centering
\includegraphics[width=0.55\linewidth]{figures/kappa_matrix.pdf}
\caption{Cohen's $\kappa$ heatmap across eleven of the twelve models (TNN omitted) on NIGHTS-100k non-tie predictions ($N=64{,}420$). Each cell is the $\kappa$ between the row and column model's binary left/right predictions on the same set of triplets. Diagonal is 1 by construction. $\kappa=1$ is perfect agreement; $\kappa=0$ is chance. Two blocks stand out: the DINO-LoRA block (within-block $\kappa>0.6$), and the two DreamSim variants that agree with each other at $\kappa=0.75$ but with LPIPS at only 0.25--0.36.}
\label{fig:kappa}
\end{figure}

\paragraph{Is $\kappa$ sufficient when labels have limited reproducibility?} One reviewer concern is that pooled $\kappa$ mixes agreement on reproducible-label triplets (X-0, ceiling 100\%) with agreement on low-reproducibility triplets (2-1, ceiling 74\%). If two models ``agree'' on a 2-1 triplet, are they agreeing on shared signal or shared noise? We address this by stratifying $\kappa$ by agreement pattern: computing $\kappa$ separately on the X-0 (unanimous), X-1 (strong, $X\geq 5$), and weak (2-1/3-1/4-1) subsets. Table~\ref{tab:kappa-strat} shows that the ranking of redundancy is preserved across all strata, with magnitudes declining in the expected way on weak-majority triplets where label noise is higher.

\begin{table}[!ht]
\caption{Cohen's $\kappa$ stratified by human-agreement pattern on the pooled NIGHTS-100k non-tie predictions. The ranking of redundancy (DreamSim-Pretrained $\leftrightarrow$ DreamSim-Scratch highest, LPIPS $\leftrightarrow$ DreamSim lowest) is identical at every stratum. On the unanimous subset, where labels have 100\% reproducibility ceiling, the DreamSim-Pretrained $\leftrightarrow$ DreamSim-Scratch agreement rises to $\kappa=0.90$, the strongest piece of evidence that retraining the ensemble on NIGHTS-100k does not alter its decision surface.}
\label{tab:kappa-strat}
\centering
\small
\resizebox{\linewidth}{!}{%
\begin{tabular}{lcccc}
\toprule
Pair & Unanimous X-0 & Strong X-1 & Weak (2-1/3-1/4-1) & Pooled \\
     & $(N=20{,}005)$ & $(N=7{,}922)$ & $(N=36{,}476)$ & $(N=64{,}420)$ \\
\midrule
DreamSim-Pretrained $\leftrightarrow$ DreamSim-Scratch & 0.90 & 0.76 & 0.68 & 0.75 \\
DreamSim-Pretrained $\leftrightarrow$ DINOv3-LoRA-Cog & 0.89 & 0.69 & 0.58 & 0.69 \\
DINOv2-LoRA-Cos $\leftrightarrow$ DINOv3-LoRA-Cog    & 0.85 & 0.64 & 0.55 & 0.67 \\
DINOv3-ZS $\leftrightarrow$ DreamSim-Pretrained         & 0.72 & 0.52 & 0.44 & 0.54 \\
LPIPS $\leftrightarrow$ DreamSim-Pretrained              & 0.52 & 0.34 & 0.27 & 0.30 \\
\bottomrule
\end{tabular}}
\end{table}

The pooled $\kappa$ in Figure~\ref{fig:kappa} and in the main body is therefore a weighted average over agreement patterns with different label reliability. The stratified view shows two things: first, the redundancy claim (DreamSim variants and DINO-LoRA pairs are structurally similar classifiers, LPIPS is independent) holds on the subset where labels are authoritative; second, the lower $\kappa$ magnitudes on weak-majority patterns partly reflect the reduced maximum $\kappa$ achievable when classifier outputs are compared against noisy labels. Both observations are consistent with the \S\ref{sec:ceiling} reproducibility-ceiling analysis: on weak-majority triplets there is less reproducible signal for any classifier (or pair of classifiers) to latch onto.

\clearpage
\noindent{\Large\bfseries D\quad Additional experiments}
\applabel{D}{app:extra}
\vspace{1em}

This section documents experiments that did not make Table~\ref{tab:main} but took substantial compute: alternative similarity heads on frozen DINOv3 CLS embeddings, a multi-teacher distillation backbone (C-RADIOv4), a vision-language-model-as-judge (Qwen3-VL), ClickMe attention harmonization, and qualitative traces that make the agreement-stratified story visible on individual triplets. The BAPPS ceiling is reported in the main body (\S\ref{sec:bapps}).

\subsection*{D.0\quad Harmonized DINOv3 (attention alignment on ClickMe)}

\citet{fel2022harmonization} propose training a vision backbone to match human attention maps collected via ClickMe~\citep{linsley2019clickme}. The hypothesis is that attention-aligned features yield similarity-aligned embeddings. We tested this on DINOv3 with the authors' pretrained ClickMe attention predictor: 5-fold CV on NIGHTS-100k non-tie (matched splits to Table~\ref{tab:main}) gives Harmonized-DINOv3 $= 64.0 \pm 0.2\%$, an 8 pp drop from DINOv3 zero-shot on the same folds (Table~\ref{tab:harm}). The near-zero across-fold variance indicates the drop is not noise.

The interpretation in the Discussion is the level-of-abstraction mismatch: ClickMe captures where humans look during categorization, not which compositional features drive perceived ensemble-image similarity. The result is consistent with attention-alignment being a useful signal for LPIPS-style pixel-distortion metrics but not for the mid/high-level discrimination NIGHTS requires.

\begin{table}[ht]
\caption{Harmonized DINOv3 on NIGHTS-100k, 5-fold CV. Zero-shot DINOv3 reproduced from Table~\ref{tab:main} for direct comparison on the same test splits.}
\label{tab:harm}
\centering
\small
\begin{tabular}{lccccccc}
\toprule
Model & F1 & F2 & F3 & F4 & F5 & Mean $\pm$ Std & $\Delta$ \\
\midrule
DINOv3-ZS-CLS-cos            & 72.4 & 72.3 & 72.1 & 71.4 & 71.7 & $72.0 \pm 0.4$ & --     \\
Harmonized DINOv3            & 63.8 & 64.1 & 63.9 & 64.4 & 63.9 & $64.0 \pm 0.2$ & $-8.0$ \\
\bottomrule
\end{tabular}
\end{table}

\subsection*{D.1\quad Alternative similarity heads on frozen CLS}

The main body establishes that cosine distance on DINO CLS tokens is the strongest zero-shot signal on NIGHTS-100k. A natural follow-up question is whether any other similarity function on the same frozen embedding improves that number. We swept ten alternatives spanning basic distances, kernel methods, non-Euclidean geometries, distribution-based interpretations, and learned metric-learning heads, all evaluated on the same 5-fold cross-validation on NIGHTS-100k non-tie (Table~\ref{tab:head-sweep}).

\begin{table}[!ht]
\caption{Alternative similarity heads on frozen DINOv3 CLS embeddings, 5-fold CV on NIGHTS-100k non-tie ($N=64{,}420$). Cosine is the baseline from Table~\ref{tab:main}. ``Learned'' heads are trained with pairwise BCE for 10 epochs; all other heads are parameter-free. No alternative beats plain cosine by more than 0.1 pp, within fold-to-fold s.d.; learned heads with more parameters do systematically worse. Retrieval-style losses did not help either (cosine-softmax-ranking and margin-triplet variants of cosine all remained within 0.1 pp of plain cosine on a fold-0 pilot).}
\label{tab:head-sweep}
\centering
\small
\begin{tabular}{llrr}
\toprule
Head & Type & Trainable params & Accuracy \\
\midrule
\textbf{Cosine} (baseline)                & parameter-free, angular & 0        & \textbf{72.01 $\pm$ 0.4} \\
Neg-dot product                           & parameter-free, inner   & 0        & 72.01 $\pm$ 0.3 \\
Arc-cosine kernel (order-1)               & Cho \& Saul 2009        & 0        & 72.11 $\pm$ 0.3 \\
Poincar\'e hyperbolic distance            & non-Euclidean           & 0        & 72.09 $\pm$ 0.3 \\
L2 Euclidean                              & parameter-free, metric  & 0        & 71.56 $\pm$ 0.3 \\
L1 Manhattan                              & parameter-free, metric  & 0        & 71.50 $\pm$ 0.2 \\
JSD on softmax(CLS)                       & distribution divergence & 0        & 71.37 $\pm$ 0.2 \\
\midrule
Mahalanobis diag (learned)                & learned metric          & 768      & 71.79 $\pm$ 0.2 \\
Mahalanobis rank-32 (learned)             & learned metric          & 24{,}576 & 70.26 $\pm$ 0.4 \\
Learned bilinear rank-32                  & learned metric          & 49{,}152 & 68.35 $\pm$ 0.3 \\
Siamese MLP on $[a, b, |a{-}b|, a{\odot}b]$ & learned similarity    & $\sim$800k & 68.39 $\pm$ 0.5 \\
\bottomrule
\end{tabular}
\end{table}

Cosine, neg-dot, arc-cosine, and Poincar\'e differ by at most 0.1 pp, less than the fold-to-fold standard deviation. L2, L1, and JSD trail cosine by $\sim$0.5--0.7 pp. Adding parameters makes performance worse: a Mahalanobis metric with 24{,}576 extra parameters loses 1.7 pp, while an 800k-parameter Siamese MLP loses 3.6 pp despite access to both embeddings. The learned heads appear to fit fold-specific label noise (\S\ref{sec:ceiling}) rather than learn a better metric. On frozen DINO CLS embeddings, no tested head beats cosine by more than 0.1 pp, and all learned heads perform worse.

\subsection*{D.2\quad Multi-teacher distillation (negative result)}

C-RADIOv4-SO400M~\citep{conwell2024large} distills DINOv3, SigLIP2, and SAM3 into a single ViT. We hypothesized that a model with stronger semantic, linguistic, and segmentation priors would produce better perceptual similarity. We applied rank-8 LoRA to qkv/proj and a cosine head on the summary token and trained for one epoch on fold 0 (approximately 28 hours at the model's required 1008$\times$1008 input resolution for its RoPE embeddings).

Epoch-1 test accuracy on fold 0 was 58.2\%, below every DINO backbone in Table~\ref{tab:main} by 13 pp or more, and below LPIPS by 7 pp. We stopped after one epoch. The model's pretraining objectives mix spatial, semantic, and linguistic signals that appear to dilute the fine-grained synthetic-image similarity signal NIGHTS requires. Training at 224$\times$224 is not possible without retraining the RoPE embeddings, and the full-resolution compute cost makes multi-fold cross-validation impractical for this backbone at our compute budget.

\subsection*{D.3\quad VLM-as-judge (negative result)}

Modern vision-language models can describe image content at the scene level. We tested whether a VLM can directly answer the 2AFC question. Using Qwen3-VL-4B-Instruct on a 100-triplet sample from fold 0, we constructed a prompt showing the reference and two candidates with randomized A/B order, and extracted the logit ratio of the A vs.\ B tokens at the first generated position. We applied position randomization to remove A/B bias.

\begin{table}[!ht]
\caption{Qwen3-VL-4B-Instruct on a 100-triplet sample of NIGHTS-100k fold 0. The stratified rows in the submitted version were incorrect and are corrected here.}
\label{tab:qwen}
\centering
\small
\begin{tabular}{lcc}
\toprule
Triplet subset & N & Accuracy \\
\midrule
Unanimous (X-0)            & 26 & 69.2\% \\
Strong-majority (X-1, $X{\geq}5$) & 13 & 84.6\% \\
Weak-majority (2-1/3-1/4-1) & 61 & 55.7\% \\
\midrule
All                        & 100 & 63.0\% (Wilson 95\% CI 53.2--71.8) \\
\bottomrule
\end{tabular}
\end{table}

Qwen3-VL reaches 63.0\% overall, below LPIPS and every zero-shot backbone in Table~\ref{tab:backbones}. On unanimous triplets it reaches 69.2\%, while every zero-shot backbone exceeds 75\%, and it is near chance on weak-majority triplets. Each agreement stratum has only 13 to 61 triplets, so only the overall figure should be read quantitatively. The model can reason about scene content but does poorly on the fine-grained similarity judgments NIGHTS asks for.

\clearpage
\noindent{\Large\bfseries D.4\quad Model confidence versus the human vote split}
\applabel{D.4}{app:confidence}
\vspace{1em}

The headline result is in \S\ref{sec:confidence}. We map each model's score to a probability $P(\text{left})$ with out-of-fold Platt scaling (fit on four folds, applied to the fifth), define model confidence as $2|P-0.5|$, and compare it to the confidence the human vote supports, $2|v-0.5|$ for left-vote fraction $v$; the 35{,}580 tied triplets never enter the fit. Expected calibration error (ECE) is reported against both the binarized majority label and the human vote fraction. Figure~\ref{fig:confcal} reports the full picture: confidence by agreement pattern against the confidence the votes themselves support, a reliability curve of calibrated $P(\text{left})$ against the observed vote fraction, and mean confidence on the 35{,}580 split-vote triplets. Figure~\ref{fig:confcum} repeats the measurement cumulatively, starting from the published 20k and relaxing the agreement filter one vote pattern at a time out to the full 100k. The confidence the votes support falls from 1.000 to 0.414 across that sweep; DreamSim-Pretrained falls from 0.790 to 0.563, and the two curves cross near 70{,}000 triplets.

\begin{table}[ht]
\caption{Mean model confidence $2|P(\text{left})-0.5|$ by agreement pattern after out-of-fold Platt scaling, against the confidence the vote split supports. X-0 pools the 2-0 to 10-0 patterns (the one 1-0 and two 11-0 triplets are excluded).}
\label{tab:confidence}
\centering
\small
\begin{tabular}{lrcccc}
\toprule
Agreement & $N$ & Vote-supported & DreamSim-Pretrained & DINOv1-ZS & LPIPS \\
\midrule
50/50 (tie) & 35{,}580 & 0.000 & 0.484 & 0.383 & 0.251 \\
2-1         & 18{,}454 & 0.333 & 0.492 & 0.386 & 0.250 \\
4-1         & 6{,}841  & 0.600 & 0.544 & 0.412 & 0.260 \\
X-0         & 20{,}016 & 1.000 & 0.790 & 0.580 & 0.324 \\
\bottomrule
\end{tabular}
\end{table}

\begin{figure}[ht]
\centering
\includegraphics[width=\linewidth]{figures/confidence_calibration.pdf}
\caption{Model confidence against the human vote split on NIGHTS-100k. (A) confidence by agreement pattern; black is the confidence the observed vote split supports. (B) reliability against the observed vote fraction; legend gives ECE against the vote fraction and against the binarized majority label. (C) mean confidence on the 35{,}580 triplets where humans split 50/50. Gray bars are LoRA-trained heads reporting their native sigmoid.}
\label{fig:confcal}
\end{figure}

\begin{figure}[ht]
\centering
\includegraphics[width=\linewidth]{figures/confidence_by_agreement.pdf}
\caption{The same measurement on the agreement split used throughout the paper. (A) the 17 raw vote patterns on the majority-vote-fraction axis of Figure~\ref{fig:agreement}; the unanimous block is the published 20k. (B) cumulative corpus, relaxing the filter one pattern at a time from the published 20k to the full 100k.}
\label{fig:confcum}
\end{figure}

\clearpage
\noindent{\Large\bfseries D.5\quad Per-axis stimulus manipulation results}
\applabel{D.5}{app:axes_detail}
\vspace{1em}

These are the per-axis results of the Figure-6-style replica described in \S\ref{sec:shape_axes}. We report them here, and not in the main body, because NIGHTS provides no human judgments on manipulated stimuli, and we collected none: the orientation, color, shading and blur axes therefore measure model invariance only and cannot be set against human behavior. The foreground/background axis is treated separately in the main text, where IN-9 and ObjectNet supply an external anchor.

\begin{figure}[ht]
\centering
\includegraphics[width=0.88\linewidth]{figures/fig6_ablation_heatmap.pdf}
\caption{DreamSim-style Figure-6 ablation replica. Accuracy drop (pp) relative to clean for each backbone $\times$ ablation. Bluer is more invariant; redder is more fragile. DINOv3 zero-shot is the most invariant overall; DreamSim-Pretrained and our own DINOv2+LoRA are among the most fragile despite the best clean accuracy. 2000 unanimous NIGHTS triplets.}
\label{fig:fig6_replica}
\end{figure}

\begin{enumerate}[leftmargin=*,itemsep=1pt,topsep=1pt]
\item \textbf{DINOv3 was the most invariant zero-shot backbone} on every axis except background removal. Rotation drop was $\leq 0.7$ pp, grayscale drop was only 0.3 pp, mild blur actually \emph{helped} (+1.25 pp at $\sigma=4$), and the foreground-only drop was 1.1 pp. Its clean accuracy of 85.4\% was below DINOv1's 88.5\% and DreamSim-Pretrained's 93.8\%, but its accuracy changed less under the tested manipulations.
\item \textbf{DreamSim-Pretrained traded invariance for clean-image accuracy.} It had the best clean score in the table ($93.8\%$) but the worst or near-worst drop on rotation ($-4.3$), color ($-5.7$), heavy blur ($-9.6$), foreground ($-9.05$), and background ($-14.5$). Margin-hinge LoRA training narrows the feature geometry around the triplet distribution, which produces the best in-distribution score at the cost of cross-manipulation stability.
\end{enumerate}

Both observations are statements about models. Neither supports a claim about the difference between model error and human error on these stimuli, and the revision no longer makes one.

\begin{table}[!ht]
\caption{Zero-shot DINOv2 with cosine on NIGHTS-100k (5-fold mean $\pm$ s.d.). Accuracy falls monotonically with size. The 70.1\% reference is the split-half agreement of the NIGHTS votes, the probability that the two halves of a triplet's votes choose the same majority; it is a conservative companion to the $\approx$89\% fresh-panel ceiling of \S\ref{sec:ceiling}.}
\label{tab:scale}
\centering
\small
\begin{tabular}{lcccc}
\toprule
DINOv2 & Params & NIGHTS-100k & vs.\ 70.1\% & Unanimous X-0 \\
\midrule
ViT-S/14 & 21M  & $72.2 \pm 0.5$ & $+2.1$ & 85.8 \\
ViT-B/14 & 86M  & $70.6 \pm 0.4$ & $+0.5$ & 83.9 \\
ViT-L/14 & 300M & $69.0 \pm 0.2$ & $-1.1$ & 81.8 \\
ViT-g/14 & 1.1B & $68.4 \pm 0.3$ & $-1.7$ & 80.5 \\
\bottomrule
\end{tabular}
\end{table}

\clearpage
\noindent{\Large\bfseries D.6\quad Zero-shot backbone sweep}
\applabel{D.6}{app:backbones}
\vspace{1em}

Every backbone below is evaluated with plain cosine on the post-norm CLS (or pooled) embedding at 224$\times$224, with no training, on the same five folds as Table~\ref{tab:main}. The DINOv2 size sweep is repeated in Table~\ref{tab:scale}.

\begin{table}[ht]
\caption{Zero-shot cosine accuracy on NIGHTS-100k non-tie (5-fold mean $\pm$ s.d., $N=64{,}420$) and on the unanimous X-0 subset.}
\label{tab:backbones}
\centering
\small
\begin{tabular}{llcc}
\toprule
Backbone & Family & NIGHTS-100k & Unanimous X-0 \\
\midrule
LPIPS-AlexNet          & CNN, supervised       & $65.3 \pm 0.2$ & 75.9 \\
ResNet-152             & CNN, supervised       & $67.1 \pm 0.4$ & 78.9 \\
CLIP ViT-B/16          & ViT, language-image   & $68.0 \pm 0.3$ & 80.0 \\
ResNet-50              & CNN, supervised       & $69.7 \pm 0.3$ & 82.4 \\
BiT-M R101x3           & CNN, supervised       & $70.2 \pm 0.3$ & 83.5 \\
BagNet-33              & CNN, bio-inspired     & $70.3 \pm 0.5$ & 83.6 \\
ConvNeXt-Large         & CNN, supervised       & $71.4 \pm 0.3$ & 85.4 \\
CORnet-S               & CNN, bio-inspired     & $71.7 \pm 0.2$ & 85.7 \\
VOneResNet-50          & CNN, bio-inspired     & $71.8 \pm 0.6$ & 85.7 \\
Supervised ViT-B/16    & ViT, supervised       & $72.4 \pm 0.2$ & 86.5 \\
\midrule
DINOv2 ViT-g/14        & ViT, self-supervised  & $68.4 \pm 0.3$ & 80.5 \\
DINOv2 ViT-L/14        & ViT, self-supervised  & $69.0 \pm 0.2$ & 81.8 \\
DINOv2 ViT-B/14        & ViT, self-supervised  & $70.6 \pm 0.4$ & 83.9 \\
DINOv3 ViT-B/16        & ViT, self-supervised  & $72.0 \pm 0.4$ & 85.7 \\
DINOv2 ViT-S/14        & ViT, self-supervised  & $72.2 \pm 0.5$ & 85.8 \\
DINOv1 ViT-B/16        & ViT, self-supervised  & $74.1 \pm 0.5$ & 88.6 \\
\midrule
DreamSim ensemble, untrained & DINO+CLIP+OpenCLIP & $74.5 \pm 0.4$ & 89.2 \\
\bottomrule
\end{tabular}
\end{table}

\noindent{\Large\bfseries D.7\quad Color control for the fine-tuning gain}
\applabel{D.7}{app:colorcontrol}
\vspace{1em}

To test whether the 4.2 pp that training adds over the untrained DreamSim ensemble (\S\ref{sec:gain}) comes from low-level color, we computed for every one of the 64{,}420 non-tie triplets the $\chi^2$ distance between the RGB histograms of each candidate and the reference, and split the triplets by whether that gap favors the human-majority candidate. The training gain was $+3.9$ pp where color could not discriminate and $+4.9$ pp where it could (95\% CI $\pm 0.7$ pp), so the gain is nearly flat across the split. A logistic regression on twelve pixel statistics alone reached 62.3\%, below every zero-shot backbone in Table~\ref{tab:backbones}.

\clearpage
\noindent{\Large\bfseries E\quad Training details}
\applabel{E}{app:training}
\vspace{1em}

This section documents the full training recipe for reproducibility. All trained models in Table~\ref{tab:main} use the same pipeline except where explicitly noted; DreamSim-Scratch uses the authors' published training code verbatim and is documented separately at the end.

\paragraph{Preprocessing.} Images are decoded from the NIGHTS release JPEGs, resized to $224\times 224$ with bicubic interpolation, and center-cropped. Normalization follows the backbone: ImageNet statistics (mean $[0.485, 0.456, 0.406]$, std $[0.229, 0.224, 0.225]$) for DINOv1, DINOv2, and DINOv3; $0.5/0.5$ for C-RADIOv4; the processor-supplied normalization for DreamSim and Qwen3-VL. No data augmentation is used at training time: the triplet task is orientation-sensitive, and the reference/candidate pairing is the supervision signal.

\paragraph{LoRA configuration.} LoRA adapters have rank 8, $\alpha = 16$, and dropout 0.05, and are inserted on the \texttt{qkv} and \texttt{proj} linear layers of every transformer block. The backbone's final LayerNorm is also unfrozen. For DINOv2 ViT-B/14 with 12 blocks this yields 442{,}368 LoRA parameters plus 1{,}536 LayerNorm parameters, which is approximately 0.5\% of the 86M backbone. The DreamSim ensemble retraining uses rank 16 on the same modules to match the authors' released configuration.

\paragraph{Heads.} Zero-shot models use cosine distance on the post-norm CLS token at $224\times224$ input. The Cog head has 197k parameters: frozen cosine alignment, GCM context, and Tversky contrast with learnable $\alpha, \beta$. The Tversky Neural Network of \citet{doumbouya2025tversky} is matched at 197k parameters, with product intersection and subtract-match difference functions and 256 features.

\paragraph{Optimizer.} Schedule-Free AdamW \citep{defazio2024schedulefree} with $\beta_1 = 0.9$, $\beta_2 = 0.999$, weight decay $10^{-2}$, 100 linear warm-up steps, and a constant learning rate thereafter (no cosine schedule). Separate learning rates are used: $10^{-4}$ for LoRA and layer-norm parameters, $2\times 10^{-4}$ for the head. Gradients are clipped to a global norm of 1.0 before each optimizer step.

\paragraph{Batch and compute.} Batch size is 32 triplets per step (each triplet contributes three images to the forward pass). Training runs for 10 epochs over the fold's training split with batch size 32. Compute is a shared cluster of NVIDIA RTX A6000 GPUs (48\,GB). A full 5-fold run of DINOv2-LoRA-Cog takes approximately 3--5 hours when five GPUs run concurrently, one per fold.

\paragraph{Checkpoint selection.} A held-out validation split of 10\% of each fold's training data selects the best-of-10-epochs checkpoint by validation accuracy; the selected checkpoint is evaluated once on the fold's test split, and that single number appears in Table~\ref{tab:perfold}. Per-triplet test-set predictions are persisted for every (model, fold) pair so that the Cohen's $\kappa$ matrix, McNemar tests, and error-distribution analyses can be reproduced without re-running training.

\paragraph{DreamSim retraining from scratch.} We run the authors' published training code unchanged with the released ensemble configuration: DINO+CLIP+OpenCLIP, rank-16 LoRA, margin-0.05 hinge loss, batch size 16, 6 epochs. The release filters training triplets to those with at least 6 votes, so the scratch model sees only strong-majority triplets during training; we evaluate on the same full non-tie test fold as every other model in Table~\ref{tab:main}. This mismatch is the point of the comparison, not an oversight: it isolates what the authors' training recipe would produce when the 100k distribution is available.

\clearpage
\noindent{\Large\bfseries F\quad Shape-only class separability and retrieval}
\applabel{F}{app:shape_only}
\vspace{1em}

\begin{table}[ht]
\caption{Shape-only class separability. Leave-one-out nearest-prototype accuracy on the Geirhos filled-silhouette and edge stimuli. Chance is 6.25\% (16 classes). Referenced from \S\ref{sec:shape_axes}.}
\label{tab:shape_only}
\centering
\small
\begin{tabular}{lcccc}
\toprule
Backbone & Silh 1-NN & Silh proto & Edges 1-NN & Edges proto \\
\midrule
ResNet-50        & 77.5 & 81.3 & 63.1 & 67.5 \\
DINOv1           & 71.3 & 78.1 & 66.3 & 68.1 \\
CLIP-ViT-B/16    & 73.1 & 78.8 & 79.4 & 81.3 \\
DINOv2           & 81.3 & 86.3 & 71.3 & 80.0 \\
DINOv3           & 82.5 & 87.5 & 86.9 & 91.3 \\
ConvNeXt-Large   & 84.4 & 87.5 & 90.0 & 93.8 \\
\bottomrule
\end{tabular}
\end{table}

\begin{figure}[ht]
\centering
\includegraphics[width=\linewidth]{figures/silhouette_retrieval.png}
\caption{Shape-only nearest-neighbor retrieval. For four query silhouettes (bicycle, cat, elephant, keyboard), each backbone's top-3 nearest silhouettes by cosine distance in CLS space. Green outline: match; red: mismatch. Every backbone, including ResNet-50 with its near-zero cue-conflict shape-bias, retrieves predominantly within-class; shape representation is present in every backbone we tested.}
\label{fig:silh_retrieval}
\end{figure}

\begin{figure}[ht]
\centering
\includegraphics[width=0.8\linewidth]{figures/cost_of_accuracy.pdf}
\caption{The cost of clean-image accuracy. Each point is a backbone; vertical axis is the mean magnitude of accuracy drop across the 11 Fig-6 ablation conditions. Zero-shot backbones (circles) cluster in the low-fragility region; LoRA-trained models and DreamSim-Pretrained (diamonds) shift up and right. Dashed line is a least-squares trend.}
\label{fig:cost_acc}
\end{figure}

\paragraph{Tie-based bias audit.} On the 7,116-triplet tie subset held out per fold, LoRA-trained models predict left with frequency 49.5--51.1\% (no positional bias). Zero-shot DINOs show 49.6--50.8\% and DreamSim 51\%, also near chance. The one model with tie bias outside that range is Qwen3-VL-4B (55\% left, confidence 0.30), consistent with known VLM positional preferences. LoRA adaptation on NIGHTS-100k does not induce spurious tie preferences.

\clearpage
\noindent{\Large\bfseries G\quad Full 2AFC-vs-JND axis table}
\applabel{G}{app:axes}
\vspace{1em}

\begin{table}[ht]
\caption{Per-backbone 2AFC accuracy (Fig-6 clean, 2000 unanimous triplets) and JND AUC (1{,}216 reference--candidate pairs from 608 triplets, same/different). Extended form of Figure~\ref{fig:tasks_decouple}. Among the zero-shot backbones the two rankings are only weakly related (Spearman $\rho = 0.37$, $p = 0.19$); the two models trained on NIGHTS 2AFC lead on both. A single AUC carries a 95\% bootstrap interval of about $\pm 0.04$.}
\label{tab:axes}
\centering
\small
\begin{tabular}{lccl}
\toprule
Backbone & 2AFC acc (\%) & JND AUC & Family \\
\midrule
DINOv1              & 88.45 & 0.674 & DINO self-distill \\
DINOv2              & 84.65 & 0.677 & DINO self-distill \\
DINOv3              & 85.35 & 0.697 & DINO self-distill \\
CLIP-ViT-B/16       & 80.60 & 0.661 & image-text \\
OpenCLIP-ViT-B/16   & 83.95 & 0.696 & image-text \\
ViT-B/16 & 86.90 & 0.670 & label-supervised ViT \\
ConvNeXt-Large      & 85.40 & 0.648 & modern CNN \\
ResNet-50           & 84.05 & 0.647 & classical CNN \\
ResNet-152          & 83.45 & 0.653 & classical CNN \\
BiT-M R101x3        & 83.50 & 0.634 & JFT-scale CNN \\
CORnet-S            & 86.60 & 0.675 & biologically-inspired \\
BagNet-33           & 84.20 & 0.651 & biologically-inspired \\
VOneResNet50        & 85.55 & 0.674 & biologically-inspired \\
DreamSim-ZS     & 89.00 & 0.676 & concatenated ensemble, ZS \\
DINOv2 + LoRA       & 89.30 & 0.720 & LoRA-trained on NIGHTS \\
DreamSim-Pretrained & 93.75 & 0.727 & LoRA-trained ensemble \\
\bottomrule
\end{tabular}
\end{table}

\clearpage
\noindent{\Large\bfseries H\quad BAPPS-vs-NIGHTS quadrant view}
\applabel{H}{app:crossover}
\vspace{1em}

\begin{figure}[ht]
\centering
\includegraphics[width=0.45\linewidth]{figures/bapps_vs_nights_crossover.pdf}
\caption{Quadrant view of Figure~\ref{fig:crossover}. Dashed lines mark the 69.0\% (BAPPS) and 70.1\% (NIGHTS) reference levels of Table~\ref{tab:bapps}. The top-right quadrant lies above both. Two models fall inside it: DreamSim-Pretrained, after supervised training on the NIGHTS-20k subset, and zero-shot DINOv1, with no training.}
\label{fig:crossover_quadrant}
\end{figure}

\begin{figure}[ht]
\centering
\includegraphics[width=0.72\linewidth]{figures/bapps_vs_nights_spectrum.pdf}
\caption{Perceptual similarity is a spectrum. Each model's accuracy at low-level BAPPS (left) connects to its accuracy at mid-to-high-level NIGHTS-100k (right). Crossing lines show rank inversion across the spectrum: LPIPS is a BAPPS specialist that trails on NIGHTS; DINOv2/v3 invert the pattern. Two models score above both reference levels: DreamSim after supervised NIGHTS-20k training, and zero-shot DINOv1 with no training.}
\label{fig:crossover}
\end{figure}

\begin{table}[ht]
\caption{Split-half agreement for BAPPS 2AFC validation (5 votes per triplet): the probability that the two halves of a triplet's votes choose the same majority, a conservative companion to the fresh-panel repeatability of \S\ref{sec:ceiling} (68.3\% for 3-2, 94.2\% for 4-1). The lower rows give per-category 2AFC accuracy against the binary majority label; their unweighted mean is the BAPPS column of Table~\ref{tab:bapps}.}
\label{tab:bapps-ceiling}
\centering
\small
\begin{tabular}{lcccccc}
\toprule
Category & CNN & Traditional & Super-res & Color & Frame-int & Deblur \\
\midrule
N        & 4{,}720 & 4{,}720 & 10{,}856 & 4{,}720 & 1{,}888 & 9{,}440 \\
Ceiling  & 83.2\% & 79.1\% & 70.0\% & 64.1\% & 62.8\% & 59.6\% \\
\midrule
LPIPS-AlexNet           & 86.7 & 77.7 & 76.2 & 66.3 & 66.9 & 63.8 \\
DINOv1-ZS               & 85.2 & 80.0 & 73.2 & 62.0 & 66.1 & 62.4 \\
DINOv2-ZS               & 81.0 & 75.4 & 63.7 & 58.1 & 66.6 & 59.1 \\
DINOv3-ZS               & 83.2 & 77.2 & 65.8 & 59.8 & 64.4 & 60.2 \\
DreamSim-Pretrained & 86.8 & 82.0 & 75.0 & 67.7 & 67.5 & 62.6 \\
\bottomrule
\end{tabular}\\[3pt]
\textbf{All BAPPS: N=36{,}344, split-half agreement = 69.0\%.}
\end{table}

\begin{table}[ht]
\caption{Accuracy on ImageNet-prompted and non-ImageNet-prompted subsets of NIGHTS-100k non-tie. $\Delta$ is ImageNet minus non-ImageNet.}
\label{tab:imnet}
\centering
\small
\begin{tabular}{lccc}
\toprule
Model & ImageNet & Non-ImageNet & $\Delta$ \\
\midrule
LPIPS-AlexNet-Pretrained       & 64.6 & 66.1 & $-1.5$ \\
DINOv1-ZS-CLS-cos              & 74.5 & 73.6 & $+0.9$ \\
DINOv3-ZS-CLS-cos              & 73.0 & 70.7 & $+2.3$ \\
DINOv2-LoRA-Cos                & 78.7 & 77.9 & $+0.7$ \\
DINOv3-LoRA-Cog             & 79.4 & 77.9 & $+1.5$ \\
DreamSim-Pretrained        & 81.0 & 80.3 & $+0.8$ \\
DreamSim-Scratch           & 79.0 & 78.3 & $+0.7$ \\
\bottomrule
\end{tabular}
\end{table}

\begin{figure}[ht]
\centering
\includegraphics[width=0.82\linewidth]{figures/error_distribution.pdf}
\caption{Left: share of each model's NIGHTS-100k errors, stratified by the human-agreement pattern of the underlying triplet (weak 2-1/3-1/4-1, strong X-1 with $X\geq 5$, unanimous X-0). Right: fraction of each model's wrong predictions where at least one human voter agreed with the model. The best models are almost never alone against the humans (DreamSim 96.6\%, DINOv3-LoRA 92.0\%); their errors live where humans themselves were split.}
\label{fig:errdist}
\end{figure}



\clearpage
\noindent{\Large\bfseries I\quad Extended \S\ref{sec:shape_axes} figures: shape bias, hGRU collapse, LoRA trajectory, 2AFC vs.\ JND}
\applabel{I}{app:shape_axes_extra}
\vspace{1em}

These five figures were referenced from \S\ref{sec:shape_axes} but deferred here for length.

\begin{figure}[ht]
\centering
\includegraphics[width=0.68\linewidth]{figures/shape_bias_cue_conflict.pdf}
\caption{Cue-conflict shape bias on the Geirhos 2018 stimuli. DINOv1 preferences are texture-dominated (10.3\%); DINOv2 (23.2\%) and DINOv3 (37.3\%) are more than three times more shape-biased. ResNet-50 is the most texture-dominant (3.75\%); ConvNeXt-Large (19.75\%) and CLIP-ViT-B/16 (29.83\%) sit between. Zero-shot, 1200 cue-conflict images, 16 classes. Humans score $\sim$95\% on the same stimuli \citep{geirhos2018imagenettrained}.}
\label{fig:shape_bias_cc}
\end{figure}

\begin{figure}[ht]
\centering
\includegraphics[width=0.78\linewidth]{figures/hgru_spatial_sweep.pdf}
\caption{Radial-mask sensitivity sweep. Accuracy on 2000 unanimous triplets as a function of radial mask radius, keeping only the center disc (fg, left) or only the surround ring (bg, right). DINO backbones transition smoothly; ResNet-50+hGRU collapses to chance on any foreground radius in $[0.1, 0.5]$ and any background radius below $0.7$, revealing that the lateral-recurrence stage learns a globally coupled representation that fails under any spatial mask.}
\label{fig:hgru_sweep}
\end{figure}

\begin{figure}[ht]
\centering
% figures/attention_zs_vs_lora.pdf is a 10-triplet strip (445.3 pt per triplet); show toaster, hotdog and spoonbill, the three the caption discusses
\includegraphics[width=0.92\linewidth,height=0.47\textheight,keepaspectratio,trim=0 3566 0 0,clip]{figures/attention_zs_vs_lora.pdf}\\[2pt]
\includegraphics[width=0.92\linewidth,height=0.235\textheight,keepaspectratio,trim=0 0 0 4008,clip]{figures/attention_zs_vs_lora.pdf}
\caption{Three high-vote NIGHTS unanimous triplets where DINOv2 + LoRA flips the majority-choice decision. Green border = human-majority candidate; red border = LoRA's (wrong) choice. Each panel is overlaid with last-layer CLS attention and annotated with the $\cos(\mathrm{ref}, \cdot)$ values the models computed. The examples share two signatures. (i) LoRA's flips occur at $\cos$ margins as small as 0.005 (spoonbill: $0.905$ vs $0.910$), exactly the voter-noise-scale differences the reproducibility analysis of Section~\ref{sec:ceiling} warned about; (ii) attention geometry is nearly identical between ZS and LoRA, so the shortcut is encoded in the feature representation (sim(full, fg-only) drops consistently for LoRA) and not in where the model looks. Pairwise similarities computed with DINOv2 + rank-8 LoRA + cosine at 224$\times$224.}
\label{fig:attn_zs_vs_lora}
\end{figure}

\begin{figure}[ht]
\centering
\includegraphics[width=0.78\linewidth]{figures/lora_strength_sweep.pdf}
\caption{Controlled LoRA-strength sweep on a single backbone. DINOv2 + rank-8 LoRA trained for 0--8 epochs on NIGHTS-20k unanimous, each checkpoint evaluated on the 2000-triplet Fig-6 ablation suite. Gray dots: 16-backbone context from Figure~\ref{fig:cost_acc}. Phase 1 (epoch 0\,$\to$\,1) mirrors the cross-backbone invariance tax slope. Phase 2 (epochs 1--8) keeps climbing the fragility axis while clean accuracy reverts toward the zero-shot baseline. Within-backbone evidence that the invariance tax is causal, not a family-level confound.}
\label{fig:lora_strength}
\end{figure}

\begin{figure}[ht]
\centering
\includegraphics[width=0.78\linewidth]{figures/tasks_decouple_scatter_v2.pdf}
\caption{2AFC accuracy and JND AUC per backbone. Horizontal axis: NIGHTS 2AFC accuracy on the 2000 unanimous-triplet subset; vertical axis: JND AUC on the 1{,}216 reference--candidate pairs of the 608 JND triplets. The green arrows trace DINOv1 $\to$ v2 $\to$ v3: v1 leads on 2AFC by 3.1--3.8 pp, but the three lie within 0.03 AUC of each other on JND, inside the $\pm 0.04$ sampling interval. The two models trained on NIGHTS 2AFC (upper right) lead on both tasks.}
\label{fig:tasks_decouple}
\end{figure}


\clearpage
\noindent{\Large\bfseries J\quad Qualitative triplet traces}
\applabel{J}{app:traces}
\vspace{1em}

The analyses in the main body are aggregate. This section shows the individual stimuli behind them, drawn unmodified from the released corpus. Figure~\ref{fig:traces_agreement} illustrates what the agreement patterns of \S\ref{sec:ceiling} look like on single triplets, one per pattern from 2-1 to 7-0. A 2-1 triplet carries a label that a fresh panel of three annotators would overturn about a quarter of the time; a 7-0 triplet carries one that reproduces every time. Both contribute equally to a headline accuracy.

\begin{figure}[!ht]
\centering
\includegraphics[width=\linewidth]{figures/triplet_traces_by_agreement_v2.png}
\caption{One NIGHTS-2AFC triplet per vote pattern, from 2-1 to 7-0: the reference, the candidate the human majority chose (green), and the alternative (red), each with its vote count. Prompts and triplet ids are given above each triplet.}
\label{fig:traces_agreement}
\end{figure}

Figure~\ref{fig:traces_hard} shows the complementary case: unanimous triplets with nine or ten votes, where the label is as stable as NIGHTS labels get, and where DreamSim-Pretrained still selects the other candidate. It does so on 20 of the 3{,}126 such triplets. These are the errors the reproducibility ceiling does not explain, and they are 3.4\% of that model's disagreements with the majority label.

\begin{figure}[!ht]
\centering
\includegraphics[width=\linewidth]{figures/unanimous_hard_examples_v2.png}
\caption{Unanimous triplets with nine or ten votes that DreamSim-Pretrained, the strongest model in Table~\ref{tab:main}, gets wrong: it picks the candidate no voter chose (red). The eight with the most votes are shown; vote pattern, prompt, and triplet id are given above each triplet.}
\label{fig:traces_hard}
\end{figure}

Figure~\ref{fig:traces_jnd} shows stimuli carrying both NIGHTS labels. All 608 shared triplets come from the X-0 unanimous tier, so the 2AFC/JND comparison of \S\ref{sec:discussion} is a comparison on the easy subset and not the hard one. JND majorities judged both candidates different from the reference in 466 triplets (77\%). In 137 triplets exactly one candidate was judged the same, the chosen one in 103 and the rejected one in 34, and in the remaining 5 both were. Comparing the two JND vote fractions within a triplet, voters judged the chosen candidate less different than the rejected one on 42.3\% of triplets, could not separate the two on 39.6\%, and judged the chosen one more different on 18.1\%.

\begin{figure}[!ht]
\centering
\includegraphics[width=\linewidth]{figures/jnd_2afc_examples_v2.png}
\caption{Triplets carrying both human labels, two for each of the three cases. The label at the top of each candidate gives its 2AFC votes; the label at the bottom gives the JND majority verdict and its vote count.}
\label{fig:traces_jnd}
\end{figure}

\clearpage

\noindent{\Large\bfseries K\quad Where human voters disagree: splits within and conflicts across the two NIGHTS tasks}
\applabel{K}{app:disagree}
\vspace{1em}

NIGHTS carries two kinds of human disagreement, and a majority label shows neither. The first is within a task: voters split on the same item. The second is across tasks: the 2AFC label and the JND label of the same image pair point in opposite directions. This section reports both from the released votes.

\paragraph{Splits within 2AFC.} Of the 64{,}420 non-tie triplets, 36{,}476 (57\%) are weak majorities (2-1, 3-1 or 4-1) and 7{,}922 have a single dissent among five or more voters; a further 35{,}580 triplets are ties. Because voting stopped at the first dissent, a split triplet usually has only three or four voters, and a fresh panel of the same size reproduces a 2-1 majority only 74\% of the time (\S\ref{sec:ceiling}). The DreamSim paper reports the split only indirectly, as the roughly 20\% of triplets that survived its unanimity filter, and the public release contains only the unanimous triplets.

\paragraph{Splits within JND.} JND voters judged each reference--candidate pair ``same'' or ``different'': three voters for 1{,}184 of the 1{,}216 pairs and six for the other 32. Of the pairs, 643 (52.9\%) were unanimous and 573 (47.1\%) were split; none was tied. Applying the fresh-panel resampling of \S\ref{sec:ceiling} to the JND votes, a split JND label is reproduced 74.7\% of the time and the average JND label 88.1\%, the same reliability profile as 2AFC. The DreamSim paper takes the majority of the three votes as the label and reports no split rate. It describes keeping triplets in which humans found exactly one candidate identical; in the release that is 137 of the 608 triplets (23\%).

\paragraph{Conflicts across the two tasks.} Table~\ref{tab:disagree} cross-tabulates the two labels on the 1{,}216 pairs. A pair is consistent when the 2AFC winner is judged ``same'' or the loser ``different,'' and in conflict when the winner is judged ``different'' or the loser ``same.'' The conflicts sit almost entirely in one place: unanimous 2AFC winners that JND voters judged different from the reference. This holds whether the JND vote was split (249 winners) or unanimous (251 winners), so it is not produced by noise in three-vote JND labels. At the triplet level, JND voters judged neither candidate ``same'' in 466 of the 608 triplets (77\%): the unanimous 2AFC vote chose between two images that JND voters saw as different from the reference. Of the 137 triplets in which exactly one candidate was judged ``same,'' that candidate was the 2AFC loser in 34 (25\%). The share of winners judged different does not fall as 2AFC agreement rises: 81\% at 6-0 and 7-0, 82\% at 8-0, 84\% at 9-0, and 91\% at 10-0 (22 triplets). Figure~\ref{fig:traces_jnd} shows examples.

\begin{table}[ht]
\caption{The two human labels on the 1{,}216 reference--candidate pairs of the 608 NIGHTS-JND triplets, all unanimous in 2AFC. Columns give the JND vote, unanimous or split, grouped by its majority. Bold cells point against the 2AFC choice (chosen judged ``different,'' rejected judged ``same''); every candidate is a visibly distorted image, so the first kind is expected, and the informative comparison is between the two rows. The last column is each row's share pointing against the 2AFC choice.}
\label{tab:disagree}
\centering
\small
\begin{tabular}{lrrrrr}
\toprule
 & \multicolumn{2}{c}{JND majority ``different''} & \multicolumn{2}{c}{JND majority ``same''} & \\
\cmidrule(lr){2-3}\cmidrule(lr){4-5}
2AFC label & unanimous & split & split & unanimous & against 2AFC \\
\midrule
Winner (chosen), $n=608$ & \textbf{251} & \textbf{249} & 92 & 16 & 82.2\% \\
Loser (rejected), $n=608$ & 372 & 197 & \textbf{35} & \textbf{4} & 6.4\% \\
\midrule
All pairs, $n=1{,}216$ & 623 & 446 & 127 & 20 & 44.3\% \\
\bottomrule
\end{tabular}
\end{table}

\paragraph{What the conflict means.} 2AFC is a forced comparison with no ``neither'' option: a candidate wins by being closer to the reference than the alternative, however different both are. JND asks whether a single candidate can be told apart from the reference. A unanimous 2AFC label therefore certifies a ranking, not similarity, and the 77\% of triplets with both candidates judged different show that most of these unanimous labels rank two visibly different images. Scoring a model against the majority label of either task discards both the splits within each task and the conflict between them.

\clearpage


\input{checklist}

\end{document}
